\subsection{Linear mappings between Lebesgue spaces defined by a kernel}

Here we take K = C. Let X be a locally compact topological space equipped with a positive measure \(\mu\) . For every \(r \in [1, +\infty]\) , we denote \(\mathcal{L}^r (\mathrm{X}) = \mathcal{L}_{\mathbf{C}}^r (\mathrm{X},\mu)\) and \(\mathrm{L}^r (\mathrm{X}) = \mathrm{L}_{\mathbf{C}}^r (\mathrm{X},\mu)\) . When \(r\in [1, + \infty [\) , one identifies the dual of \(\mathrm{L}^r (\mathrm{X})\) with \(\mathrm{L}^{r^{\prime}}(\mathrm{X})\) where \(r^\prime\) is the conjugate exponent of \(r\) such that \(1 / r + 1 / r' = 1\) (INT, V, p. 61, § 5, no. 8, th. 4). We denote by \(\| f\| _r\) the norm (or semi-norm) of \(f\in \mathrm{L}^r (\mathrm{X})\) (or \(f\in \mathcal{L}^r (\mathrm{X})\) ).

Recall that a subset A of X is said to be \(\mu\) -moderate (INT, V, p. 4, § 1, n° 2) if it is contained in the union of a sequence of \(\mu\) -integrable sets, and that a function \(f\) defined on X with values in a vector space or in \(\overline{\mathbf{R}}\) is \(\mu\) -moderate if it vanishes on the complement of a \(\mu\) -moderate subset of X. If \(1 \leqslant p < +\infty\) and \(f \in \mathcal{L}^p(\mathrm{X})\) , then \(f\) is moderate (INT, V, p. 9, § 1, no. 3, cor.). If \(f\) is \(\mu\) -moderate, then \(fg\) is \(\mu\) -moderate for every function \(g\) .

Let \(p\) and \(q\) be elements of \(]1, +\infty[\) . Let X and Y be locally compact topological spaces, equipped with positive measures \(\mu\) and \(\nu\) , respectively. One equips the space \(X \times Y\) (resp. \(Y \times X\) ) with the product measure \(\mu \otimes \nu\) (resp. \(\nu \otimes \mu\) ).

We denote \(\mathcal{N}^{p,q}(\mathrm{X}\times\mathrm{Y},\mu\otimes\nu)\) , or simply \(\mathcal{N}^{p,q}(\mathrm{X}\times\mathrm{Y})\) , the set of functions k measurable with respect to \((\mu\otimes\nu)\) from X × Y into C such that there exists a real number \(C\geqslant0\) satisfying

\[
\int_ {\mathrm{X} \times \mathrm{Y}} ^ {*} | k (x, y) f (x) g (y) | d (\mu \otimes \nu) (x, y) \leqslant \mathrm{C} \| f \| _ {p} \| g \| _ {q}\tag{2}
\]

for all functions \(f \in \mathcal{L}^p(\mathrm{X})\) and \(g \in \mathcal{L}^q(\mathrm{Y})\) . We then denote by \(\| k \|_{p,q}\) the lower bound of the set of real numbers C having this property.

Remark. --- Let \(f \in \mathcal{L}^{p}(\mathrm{X})\) and \(g \in \mathcal{L}^{q}(\mathrm{Y})\) . Since \(p < +\infty\) and \(q < +\infty\) , the functions f and g are moderate, and the function defined by \((x, y) \mapsto f(x)g(y)\) is \((\mu \otimes \nu)\) -moderate (cf. INT, V, p. 92, § 8, n°3, cor. 2). Consequently, the function \((x, y) \mapsto k(x, y)f(x)g(y)\) is \((\mu \otimes \nu)\) -moderate, and in particular the upper integral of its absolute value coincides with the essential upper integral (INT, V, p. 6, § 1, n°3, prop. 7).

The set \(\mathcal{N}^{p,q}(\mathrm{X}\times \mathrm{Y})\) is a vector subspace of the space of functions from \(\mathrm{X}\times\mathrm{Y}\) into \(\mathbf{C}\) and the mapping defined by \(k\mapsto \| k\|_{p,q}\) is a semi-norm on \(\mathcal{N}^{p,q}(\mathrm{X}\times \mathrm{Y})\) .

Lemma 1. --- Let \(k \in \mathcal{N}^{p,q}(X \times Y)\) . The function \(k\) is locally \((\mu \otimes \nu)\) -integrable.

The map \(k\) is \((\mu \otimes \nu)\) -measurable by hypothesis. Let A be a compact subset of \(\mathrm{X} \times \mathrm{Y}\) . There exist compact subsets \(\mathrm{B} \subset \mathrm{X}\) and C ⊂ Y such that A ⊂ B × C. Applying (2) to the characteristic functions \(\varphi_{B}\) of B and \(\varphi_{C}\) of C, we obtain

\[
\begin{array}{l} \int_ {X \times Y} ^ {*} | k (x, y) | \varphi_ {\mathrm{A}} (x, y) d (\mu \otimes \nu) (x, y) \\ \qquad \leqslant \int_ {X \times Y} ^ {*} | k (x, y) | \varphi_ {\mathrm{B}} (x) \varphi_ {\mathrm{C}} (y) d (\mu \otimes \nu) (x, y) <   + \infty , \end{array}
\]

whence the result (INT, V, p. 41, §5, n° 1, prop. 1 and def. 1).

PROPOSITION 3. --- a) Let \(k \in \mathcal{N}^{p,q}(X \times Y)\) . For \(f \in \mathcal{L}^p(X)\) and \(g \in \mathcal{L}^q(Y)\) , the function

\[
(x, y) \mapsto k (x, y) f (x) g (y)
\]

is \((\mu \otimes \nu)\) -integrable and there exists a unique map \(u_{k}\) from \(\mathrm{L}^{p}(\mathrm{X})\) into \(\mathrm{L}^{q'}(\mathrm{Y})\) such that

\[
\langle g, u _ {k} (f) \rangle = \int_ {\mathrm{X} \times \mathrm{Y}} k (x, y) f (x) g (y) d (\mu \otimes \nu) (x, y)
\]

for all \(f \in \mathrm{L}^p(\mathrm{X})\) and \(g \in \mathrm{L}^q(\mathrm{Y})\) . The map \(u_k\) is linear and continuous, its norm is \(\leqslant \| k\|_{p,q}\) ;

\begin{enumerate}
\def\labelenumi{\alph{enumi})}
\setcounter{enumi}{1}
\item
  Let \(k \in \mathcal{N}^{p,q}(X \times Y)\) . We have \(u_k = 0\) if and only if k is locally \((\mu \otimes \nu)\) -negligible;
\item
  The map which to k associates \(u_{k}\) is a continuous linear map from the semi-normed space \(\mathcal{N}^{p,q}(\mathrm{X}\times\mathrm{Y})\) to the Banach space \(\mathcal{L}(\mathrm{L}^{p}(\mathrm{X});\mathrm{L}^{q^{\prime}}(\mathrm{Y}))\) , such that \(\|u_{k}\|\leqslant\|k\|_{p,q}\) for all \(k\in\mathcal{N}^{p,q}(\mathrm{X}\times\mathrm{Y})\) .
\end{enumerate}

Let \(k \in \mathcal{N}^{p,q}(\mathrm{X} \times \mathrm{X})\) . By definition, the function from \(\mathrm{X} \times \mathrm{Y}\) into \(\mathbf{C}\) defined by \((x,y) \mapsto k(x,y)f(x)g(y)\) is \((\mu \otimes \nu)\) -integrable when \(f \in \mathcal{L}^p(\mathrm{X})\) and \(g \in \mathcal{L}^q(\mathrm{Y})\) , and the map

\[
b _ {k} \colon (f, g) \mapsto \int_ {\mathrm{X} \times \mathrm{Y}} k (x, y) f (x) g (y) d (\mu \otimes \nu) (x, y)
\]

is a continuous bilinear map from \(\mathrm{L}^{p}(\mathrm{X})\times\mathrm{L}^{q}(\mathrm{Y})\) to C. For every \(f\in\mathrm{L}^{p}(\mathrm{X})\) , there therefore exists a unique \(h\in\mathrm{L}^{q^{\prime}}(\mathrm{Y})\) such that \(b_{k}(f,g)=\langle g,h\rangle\) for all \(g\in\mathrm{L}^{q}(\mathrm{Y})\) ; denote it \(u_{k}(f)=h\) . By (2), we have \(\|u_{k}(f)\|_{q^{\prime}}\leqslant\|k\|_{p,q}\|f\|_{p}\) ; the map \(u_{k}\) is linear and continuous with norm \(\leqslant\|k\|_{p,q}\) . Moreover, the map \(k\mapsto u_{k}\) is linear, and assertions a) and c) follow from this.

Let us prove \(b\) ). Suppose that \(k\) is locally \((\mu \otimes \nu)\) -negligible. Let \(f \in \mathcal{L}^p(\mathrm{X})\) and \(g \in \mathcal{L}^q(\mathrm{Y})\) ; the function defined on \(\mathrm{X} \times \mathrm{Y}\) by \((x, y) \mapsto k(x, y)f(x)g(y)\) is \((\mu \otimes \nu)\) -moderate. Since it is locally \((\mu \otimes \nu)\) -negligible, it is negligible (INT, V, p. 7, § 1, n°2, cor. 1). Thus, we have \(\langle u_k(f), g \rangle = b_k(f, g) = 0\) . It follows that \(u_k = 0\) .

Conversely, let \(k \in \mathcal{N}^{p,q}(X \times Y)\) be such that \(u_k = 0\) . The function k is locally \((\mu \otimes \nu)\) -integrable (Lemma 1). For all functions \(f \in \mathcal{K}(X)\) and \(g \in \mathcal{K}(Y)\) , we have

\[
\langle g, u _ {k} (f) \rangle = \int_ {\mathrm{X} \times \mathrm{Y}} f (x) g (y) k (x, y) d (\mu \otimes \nu) (x, y),
\]

therefore the measure \(k \cdot (\mu \otimes \nu)\) on X × Y is zero (cf. INT, III, p. 82, § 4, no. 1, Th. 1). This means that k is locally \((\mu \otimes \nu)\) -negligible (INT, V, p. 46, §5, no. 3, Cor. 2).

DEFINITION 1. --- Let \(k \in \mathcal{N}^{p,q}(\mathrm{X} \times \mathrm{Y})\) . We say that the map \(u_k\) is the integral operator with kernel \(k\) from \(\mathrm{L}^p(\mathrm{X})\) into \(\mathrm{L}^{q'}(\mathrm{Y})\) .

For \(k \in \mathcal{N}^{p,q}(\mathrm{X} \times \mathrm{Y})\) and \(f \in \mathcal{L}^p(\mathrm{X})\) , we will sometimes also denote by \(u_k(f)\) the image under \(u_k\) of the class of \(f\) in \(\mathrm{L}^p(\mathrm{X})\) .

PROPOSITION 4. --- Let \(k \in \mathcal{N}^{p,q}(\mathrm{X} \times \mathrm{Y})\) and \(f \in \mathcal{L}^p(\mathrm{X})\) . The mapping defined on \(\mathrm{X}\) by \(x \mapsto k(x,y)f(x)\) is \(\mu\) -integrable for every \(y\) outside a locally \(\nu\) -negligible set, and \(u_k(f)\) coincides with the class of the function

\[
y \mapsto \int_ {\mathrm{X}} k (x, y) f (x) d \mu (x)
\]

defined locally ν-almost everywhere.

Let \(Y'\) denote the set of \(y \in Y\) such that the mapping \(x \mapsto k(x, y)f(x)\) is not \(\mu\) -integrable. Let \(y \in Y\) and let C be a compact neighborhood of y, whose characteristic function we denote by \(\varphi\) . By condition (2) applied to f and \(\varphi\) , the function \((x, y) \mapsto k(x, y)f(x)\varphi(y)\) is integrable on X × Y with respect to the measure \(\mu \otimes \nu\) . By the Lebesgue--Fubini theorem (INT, V, p. 96, § 8, n°4, th. 1, a)), the function \(x \mapsto k(x, y)f(x)\varphi(y)\) is \(\mu\) -integrable for y outside a \(\nu\) -negligible set \(Y_{C}\) . Since \(Y' \cap C \subset Y_{C}\) , it follows that \(Y'\) is locally \(\nu\) -negligible in Y (INT, IV, p. 172, § 5, n°2, def. 3).

Let h denote the mapping defined locally \(\nu\) -almost everywhere on Y by \(y \mapsto \int k(x, y)f(x)d\mu\) . It is locally \(\nu\) -integrable (INT, V, loc.

cit.). Let \(g \in \mathcal{K}(\mathrm{Y})\) . We have

\[
\begin{array}{l} \langle g, u _ {k} (f) \rangle = \int_ {\mathrm{X} \times \mathrm{Y}} k (x, y) f (x) g (y) d (\mu \otimes \nu) (x, y) \\ \qquad = \int_ {\mathrm{Y}} \Bigl (\int_ {\mathrm{X}} k (x, y) f (x) g (y) d \mu (x) \Bigr) d \nu (y) \\ \qquad = \int_ {\mathrm{Y}} g (y) h (y) d \nu (y) \end{array}
\]

(INT, V, loc. cit.). We therefore have \(u_{k}(f) = h\) locally \(\nu\) -almost everywhere (cf. INT, V, p. 46, § 5, n°3, cor. 2).

PROPOSITION 5. --- The space \(\mathcal{L}^{p'}(\mathrm{X}\times\mathrm{Y})\) is contained in \(\mathcal{N}^{p,p}(\mathrm{X}\times\mathrm{Y})\) . The linear map from \(\mathcal{L}^{p'}(\mathrm{X}\times\mathrm{Y})\) into \(\mathcal{L}(\mathrm{L}^{p}(\mathrm{X});\mathrm{L}^{p'}(\mathrm{Y}))\) defined by \(k\mapsto u_{k}\) induces, by passing to the quotient, a continuous injective linear map from \(\mathrm{L}^{p'}(\mathrm{X}\times\mathrm{Y})\) into \(\mathcal{L}(\mathrm{L}^{p}(\mathrm{X});\mathrm{L}^{p'}(\mathrm{Y}))\) .

Let \(k \in \mathcal{L}^{p'}(\mathrm{X} \times \mathrm{Y})\) . The function \(k\) is measurable (INT, IV, p. 84, § 5, n° 6, th. 5). Let \(f \in \mathcal{L}^p(\mathrm{X})\) and \(g \in \mathcal{L}^p(\mathrm{Y})\) . The function \(h \colon (x,y) \mapsto f(x)g(y)\) belongs to \(\mathcal{L}^p(\mathrm{X} \times \mathrm{Y})\) and \(\|h\|_p = \|f\|_p \|g\|_p\) (INT, V, p. 95, § 8, n° 3, cor. 2). The function \(hk\) is therefore integrable on \(\mathrm{X} \times \mathrm{Y}\) and we have

\[
(\mu \otimes \nu) ^ {*} (| h k |) \leqslant \| k \| _ {p ^ {\prime}} \| h \| _ {p} = \| k \| _ {p ^ {\prime}} \| f \| _ {p} \| g \| _ {p}
\]

by Hölder's inequality (INT, IV, p. 208, § 6, n° 4, th. 2), which proves that \(k \in \mathcal{N}^{p,p}(\mathrm{X} \times \mathrm{Y})\) with \(\| k \|_{p,p} \leqslant \| k \|_{p'}\) .

Since \(p' \neq +\infty\) , every function \(k \in \mathcal{L}^{p'}(\mathrm{X} \times \mathrm{Y})\) is moderate (INT, V, p. 9, § 1, n° 3, cor.), and the last assertion then follows from prop. 3, b) and c).

For any function \(k\colon X \times Y \to \mathbf{C}\) and for every \(y \in Y\) , we denote by \(k_y\) the mapping from X into \(\mathbf{C}\) given by \(k_y(x) = k(x, y)\) , and we denote by \(k_\circ\) the mapping from Y into \(\mathcal{F}(X; \mathbf{C})\) defined by \(y \mapsto k_y\) . We shall likewise denote the mapping from Y into the space of functions defined \(\mu\) -almost everywhere on X that associates with \(y\) the class of \(k_y\) .

For any function \(k\) from \(\mathrm{X} \times \mathrm{Y}\) into \(\mathbf{C}\) , we denote

\[
\mathrm{N} _ {p ^ {\prime}, q ^ {\prime}} (k) = \left(\int_ {\mathrm{Y}} ^ {\bullet} \mathrm{N} _ {p ^ {\prime}} (k _ {y}) ^ {q ^ {\prime}} d \nu (y)\right) ^ {1 / q ^ {\prime}}.
\]

This is an element of \([0, +\infty]\) .

Lemma 2. --- For every \(k \in \mathcal{L}^{p'}(\mathrm{X} \times \mathrm{Y})\) , we have \(\mathrm{N}_{p', p'}(k) = \| k \|_{p'}\) .

Let \(k \in \mathcal{L}^{p'}(\mathrm{X} \times \mathrm{Y})\) . By INT, V, p. 96, § 8, no. 4, th. 1, a), the set \(Y'\) of elements \(y \in \mathrm{Y}\) such that the function \(|k_y|^{p'}\) is not \(\mu\) -integrable is \(\nu\) -negligible; moreover, the function defined by \(y \mapsto \mu(|k_y|^{p'})\) on \(\mathrm{Y} - \mathrm{Y}'\) and equal to zero on \(\mathrm{Y}'\) is \(\nu\) -integrable and satisfies

\[
\| k \| _ {p ^ {\prime}} ^ {p ^ {\prime}} = \int_ {\mathrm{X} \times \mathrm{Y}} | k (x, y) | ^ {p ^ {\prime}} d (\mu \otimes \nu) (x, y) = \int_ {\mathrm{Y}} \mu (| k _ {y} | ^ {p ^ {\prime}}) d \nu = \mathrm{N} _ {p ^ {\prime}, p ^ {\prime}} (k) ^ {p ^ {\prime}}.
\]

We denote \(\mathcal{L}^{p',q'}(\mathrm{X},\mathrm{Y},\mu,\nu)\) , or simply \(\mathcal{L}^{p',q'}(\mathrm{X},\mathrm{Y})\) , the complex vector space of functions \((\mu\otimes\nu)\) -measurable \(k\) from \(\mathrm{X}\times\mathrm{Y}\) into \(\mathbf{C}\) such that, for \(\nu\) -almost all \(y\in\mathrm{Y}\) , the map \(k_y\) belongs to \(\mathcal{L}^{p'}(\mathrm{X})\) , and such that the map \(k_\circ\) belongs to \(\mathcal{L}^{q'}_{\mathrm{L}^{p'}(\mathrm{X})}(\mathrm{Y},\nu)\) (INT, IV, p. 129, § 3, no. 4, def. 2). We equip \(\mathcal{L}^{p',q'}(\mathrm{X},\mathrm{Y})\) with the seminorm \(k\mapsto\mathrm{N}_{p',q'}(k)\) ; note that this then satisfies \(\mathrm{N}_{p',q'}(k)=\|k_\circ\|_{q'}\) , where \(k_\circ\) is viewed as a mapping with values in the Banach space \(\mathrm{L}^{p'}(\mathrm{X})\) .

PROPOSITION 6. --- a) The space \(\mathcal{L}^{p'}(\mathrm{X} \times \mathrm{Y})\) is contained in \(\mathcal{L}^{p',p'}(\mathrm{X},\mathrm{Y})\) and for \(k \in \mathcal{L}^{p'}(\mathrm{X} \times \mathrm{Y})\) , we have \(\mathrm{N}_{p',p'}(k) = \| k\|_{p'}\) ;

\begin{enumerate}
\def\labelenumi{\alph{enumi})}
\setcounter{enumi}{1}
\tightlist
\item
  The space \(\mathcal{L}^{p',q'}(\mathrm{X},\mathrm{Y})\) is contained in \(\mathcal{N}^{p,q}(\mathrm{X}\times \mathrm{Y})\) and for \(k\in \mathcal{L}^{p',q'}(\mathrm{X},\mathrm{Y})\) , we have \(\| k\|_{p,q}\leqslant \mathrm{N}_{p',q'}(k)\) .
\end{enumerate}

Let us prove assertion \(a\) ). Let \(k \in \mathcal{L}^{p'}(\mathrm{X} \times \mathrm{Y})\) . By INT, V, p. 96, § 8, no. 4, th. 1, \(a\) ), the function \(k_y\) belongs to \(\mathcal{L}^{p'}(\mathrm{X})\) for \(\nu\) -almost all \(y \in \mathrm{Y}\) .

Let \(\varepsilon > 0\) . There exists \(h \in \mathcal{K}(\mathrm{X} \times \mathrm{Y})\) such that \(\|k - h\|_{p'} < \varepsilon\) . The function \(h_{\circ}\) is a continuous function with compact support from Y into \(\mathrm{L}^{p'}(\mathrm{X})\) satisfying

\[
\| k _ {\circ} - h _ {\circ} \| _ {p ^ {\prime}} = \mathrm{N} _ {p ^ {\prime}, p ^ {\prime}} (k - h) = \| k - h \| _ {p ^ {\prime}} <   \varepsilon
\]

by Lemma 2. It follows that \(k_{\circ}\) belongs to \(\mathcal{L}_{\mathrm{L}^{p^{\prime}}(\mathrm{X})}^{p^{\prime}}(\mathrm{Y},\nu)\) , that is, \(k\) belongs to \(\mathcal{L}^{p^{\prime},p^{\prime}}(\mathrm{X},\mathrm{Y})\) . Lemma 2 then shows that \(\mathrm{N}_{p^{\prime},p^{\prime}}(k)=\|k\|_{p^{\prime}}\) .

Let us prove assertion \(b\) ). Let \(k \in \mathcal{L}^{p', q'}(\mathrm{X}, \mathrm{Y})\) and let \(f \in \mathcal{L}^p(\mathrm{X})\) . For every \(y \in \mathrm{Y}\) , Hölder's inequality (INT, IV, p. 208, § 6, no. 4, th. 2) applied to \(k_y\) implies

\[
\int_ {X} ^ {*} | k (x, y) f (x) | d \mu (x) \leqslant \| k _ {y} \| _ {p ^ {\prime}} \| f \| _ {p}.
\]

Let \(g \in \mathcal{L}^{q'}(\mathrm{Y})\) . The maps \(f\) and \(g\) are moderate, since \(p\) and \(q'\) are finite. Consequently, the map \((x,y) \mapsto f(x)g(y)\) is \((\mu \otimes \nu)\) -moderate (cf. INT, V, p. 92, § 8, no. 3, cor. 2), and hence the map \((x,y)\mapsto k(x,y)f(x)g(y)\) is \((\mu\otimes\nu)\) -moderate. According to INT, V, p. 93, § 8, no. 3, prop. 7, a), it then follows that

\[
\begin{array}{l} \int_ {\mathrm{X} \times \mathrm{Y}} ^ {*} | k (x, y) f (x) g (y) | d (\mu \otimes \nu) \\ = \int_ {\mathrm{Y}} ^ {*} | g (y) | \Big (\int_ {\mathrm{X}} ^ {*} | k (x, y) f (x) | d \mu (x) \Big) d \nu (y) \\ \leqslant \| f \| _ {p} \int_ {\mathrm{Y}} ^ {*} \| k _ {y} \| _ {p ^ {\prime}} | g (y) | d \nu (y) \leqslant \mathrm{N} _ {p ^ {\prime}, q ^ {\prime}} (k) \| f \| _ {p} \| g \| _ {q}, \end{array}
\]

using Hölder's inequality once again. This concludes the proof.

PROPOSITION 7. --- Let \(k \in \mathcal{L}^{p',q'}(\mathrm{X},\mathrm{Y})\) . The linear map \(u_k\) from \(\mathrm{L}^p (\mathrm{X})\) into \(\mathrm{L}^{q'}(\mathrm{Y})\) is compact.

Suppose first that \(k\) is such that \(k_{\circ} \in \mathcal{K}(\mathrm{Y}; \mathcal{L}^{p'}(\mathrm{X}))\) and denote by A the support of \(k_{\circ}\) . Let F be the subspace of \(\mathcal{K}(\mathrm{Y}, \mathrm{A}; \mathcal{L}^{p'}(\mathrm{X}, \mu))\) generated by the functions \(y \mapsto f_2(y)f_1\) , where \(f_2 \in \mathcal{K}(\mathrm{Y}, \mathrm{A})\) and \(f_1 \in \mathcal{L}^{p'}(\mathrm{X})\) . When \(m \in \mathcal{K}(\mathrm{Y}; \mathcal{L}^{p'}(\mathrm{X}))\) satisfies \(m_{\circ} \in \mathrm{F}\) , the linear map \(u_m\) is of finite rank, hence compact. For \(m_1\) and \(m_2\) in \(\mathcal{K}(\mathrm{Y}, \mathrm{A}; \mathcal{L}^{p'}(\mathrm{X}))\) we also have

\[
\begin{array}{c} \| u _ {m _ {1}} - u _ {m _ {2}} \| \leqslant \| m _ {1} - m _ {2} \| _ {p ^ {\prime}, q ^ {\prime}} = \Big (\int_ {\mathrm{Y}} \| m _ {1, y} - m _ {2, y} \| _ {p ^ {\prime}} ^ {q ^ {\prime}} d \nu (y) \Big) ^ {1 / q ^ {\prime}} \\ \leqslant \nu (\mathrm{A}) ^ {1 / q ^ {\prime}} \sup _ {y \in \mathrm{Y}} \mathrm{N} _ {p ^ {\prime}} (m _ {1, y} - m _ {2, y}). \end{array}
\]

Since the space F is dense in \(\mathcal{K}(\mathrm{Y},\mathrm{A};\mathrm{L}^{p^{\prime}}(\mathrm{X}))\) for the topology of uniform convergence on A (INT, III, p. 41, § 1, n° 1, prop. 1 and INT, III, p. 46, § 1, n° 2, prop. 5), one concludes that the linear map \(u_{k}\) is the limit of a sequence of finite-rank linear maps. It is therefore compact (cor. 1 of III, p. 4).

Consider the general case. Let \(k \in \mathcal{L}^{p', q'}(\mathrm{X}, \mathrm{Y})\) . For every \(\varepsilon > 0\) there exists a map \(k_{\varepsilon, \circ} \in \mathcal{K}(\mathrm{Y}; \mathcal{L}^{p'}(\mathrm{X}))\) such that

\[
\left\| k _ {\circ} - k _ {\varepsilon , \circ} \right\| _ {q ^ {\prime}} <   \varepsilon .
\]

The corresponding function \(k_{\varepsilon}\) : X × Y → C then satisfies

\[
\| u _ {k} - u _ {k _ {\varepsilon}} \| \leqslant \| k - k _ {\varepsilon} \| _ {p, q} \leqslant \mathrm{N} _ {p ^ {\prime}, q ^ {\prime}} (k - k _ {\varepsilon}) = \| k _ {\circ} - k _ {\varepsilon , \circ} \| _ {q ^ {\prime}} <   \varepsilon .
\]

Since \(u_{k_{\varepsilon}}\) is compact by the preceding, the map \(u_{k}\) is also compact (prop. 2 of III, p. 4).

\pandocbounded{\includegraphics[keepaspectratio]{images/part001_ec4b93bc.jpg}}

Set \(p = 1\) and suppose \(q > 1\) . There may exist mappings \(k \colon \mathrm{X} \times \mathrm{Y} \to \mathbf{C}\) such that \(k_y \in \mathcal{L}^\infty(\mathrm{X})\) for all \(y \in \mathrm{Y}\) , and such that the integral

\[
\int_ {Y} ^ {*} \| k _ {y} \| _ {\infty} ^ {q ^ {\prime}} d \nu (y)
\]

is finite (in particular, condition (2) of III, p. 27 is valid), but such that the linear mapping \(u_k\) from \(L^1(X)\) into \(L^{q'}(X)\) is not compact (exercise 2 of III, p. 119).

In particular, one deduces from the proposition the following corollary:

COROLLARY 1 (Hilbert--Schmidt). --- Let \(k \in \mathcal{L}^{2}(X \times Y)\) . The linear map \(u_{k}\) is compact from \(L^{2}(X)\) into \(L^{2}(Y)\) .

By prop. 5, one has \(k \in \mathcal{N}^{2,2}(\mathrm{X} \times \mathrm{Y})\) , so the linear map \(u_k\) is defined (prop. 3). One has \(k \in \mathcal{L}^{2,2}(\mathrm{X} \times \mathrm{Y})\) (prop. 6, \(a\) ), hence \(u_k\) is compact (prop. 7).

Remark. --- In the case p = q = 2, it is generally more convenient to express the characterization of the integral operator \(u_{k}\) using the scalar product: it is the unique map from \(\mathrm{L}^{2}(\mathrm{X})\) into \(\mathrm{L}^{2}(\mathrm{Y})\) such that

\[
\langle g \mid u _ {k} (f) \rangle = \int_ {\mathrm{X} \times \mathrm{Y}} k (x, y) f (x) \overline {{g (y)}} d (\mu \otimes \nu) (x, y)
\]

for all \(f \in \mathrm{L}^{2}(\mathrm{X})\) and \(g \in \mathrm{L}^{2}(\mathrm{Y})\) . Moreover, the adjoint of \(u_{k}\) is the linear map \(u_{k^{*}}\) , where \(k^{*} \in \mathrm{L}^{2}(\mathrm{Y} \times \mathrm{X})\) satisfies \(k^{*}(y, x) = \overline{k(x, y)}\) for almost every \((y, x) \in \mathrm{Y} \times \mathrm{X}\) . Indeed, for all \(f \in \mathrm{L}^{2}(\mathrm{X})\) and \(g \in \mathrm{L}^{2}(\mathrm{Y})\) , we have

\[
\langle g \mid u _ {k} (f) \rangle = \int_ {\mathrm{X} \times \mathrm{Y}} k (x, y) \overline {{g (y)}} f (x) d (\mu \otimes \nu) (x, y) = \langle u _ {k ^ {*}} (g) \mid f \rangle .
\]

COROLLARY 2. --- Let G be a compact topological group equipped with a Haar measure μ. Let p be a real number such that 1 \textless{} p \textless{} +∞ and let q be the conjugate exponent of p. Let \(f \in \mathcal{L}^{q}(G)\). Set k(x,y) = f(x⁻¹y) for all (x,y) ∈ G × G.

\begin{enumerate}
\def\labelenumi{\alph{enumi})}
\item
  We have \(k \in \mathcal{L}^{q,p}(\mathrm{G},\mathrm{G})\) ;
\item
  For \(\varphi \in \mathrm{L}^p (\mathrm{G})\) the convolution \(\varphi *f\) belongs to \(\mathrm{L}^p (\mathrm{G})\) and the linear map \(v_{f}\colon \varphi \mapsto \varphi *f\) of \(\mathrm{L}^p (\mathrm{G})\) into itself is continuous. It coincides with the endomorphism \(u_{k}\) . In particular, the linear map \(v_{f}\) is compact.
\end{enumerate}

Since G is compact, the function f belongs to \(\mathcal{L}^{1}(G)\) . The linear map \(v_{f}\) is therefore defined and continuous according to INT, VIII, p. 167, § 4, n° 5, prop. 13.

The map \(k\) is \((\mu \otimes \mu)\) -measurable, and we have \(\| k_y\| _q = \| f\| _q\) for all \(y\in G\) . Since the measure \(\mu\) is bounded, it follows that

\[
\int_ {\mathrm{G}} \| k _ {y} \| _ {q} ^ {p} d \mu (y) = \mu (\mathrm{G}) \| f \| _ {q} ^ {p},
\]

hence \(k \in \mathcal{L}^{q,p}(G, G)\) . The linear map \(u_{k}\) is therefore a compact linear map from \(\mathrm{L}^{p}(\mathrm{G})\) into \(\mathrm{L}^{p}(\mathrm{G})\) (prop. 7).

Let \(\varphi\in\mathcal{K}(G)\) . By INT, VIII, p. 166, § 4, no. 5, prop. 11 and prop. 4 of III, p. 29, we then have

\[
\begin{array}{c} v _ {f} (\varphi) (y) = \int_ {\mathrm{G}} \varphi (x) f (x ^ {- 1} y) d \mu (x) \\ = \int_ {\mathrm{G}} \varphi (x) k (x, y) d \mu (x) = u _ {k} (\varphi) (y) \end{array}
\]

for almost every \(y \in G\) . This implies that \(v_{f} = u_{k}\) .

\subsection{Restriction of differentiable mappings}

Let n and r be positive integers, let U be an open subset of \(R^{n}\), and let F be a Banach space. Denote by \(\mathcal{C}^{r}(\mathrm{U};\mathrm{F})\) the vector space of maps of class \(C^{r}\) from U into F, endowed with the topology of \(C^{r}\) compact convergence. Recall that this is the supremum of the topologies of \(C^{r}\)-uniform convergence on K (VAR, R2, 12.3.10, p. 56), as K ranges over the set of compact subsets of U. The space \(\mathcal{C}^{0}(\mathrm{U};\mathrm{F})\) is none other than the space \(\mathcal{C}(\mathrm{U};\mathrm{F})\) of continuous maps from U into F, equipped with the topology of compact convergence.

Lemma 3. --- Let A be the set of multi-indices \(\alpha \in \mathbf{N}^n\) such that \(|\alpha| \leqslant r\) and let u be the linear map \(f \mapsto (\partial^\alpha f)_{\alpha \in \mathrm{A}}\) from \(\mathcal{C}^r(\mathrm{U};\mathrm{F})\) into \(\mathcal{C}(\mathrm{U};\mathrm{F})^{\mathrm{A}}\) .

\begin{enumerate}
\def\labelenumi{\alph{enumi})}
\item
  The map u is injective, continuous, strict, and its image is closed.
\item
  The topological vector space \(\mathcal{C}^{r}(\mathrm{U};\mathrm{F})\) is complete.
\end{enumerate}

The map \(u\) is linear and injective. It is continuous and strict by definition of the topology of \(\mathcal{C}^r (\mathrm{U};\mathrm{F})\) .

Let \((g_{\alpha})_{\alpha\in\mathrm{A}}\) be a point of \(\mathcal{C}(\mathrm{U};\mathrm{F})^{\mathrm{A}}\) adherent to the image of u. There exists a filter F on \(\mathcal{C}^{r}(\mathrm{U};\mathrm{F})\) such that one has \(g_{\alpha}=\lim_{f,\mathcal{F}}\partial^{\alpha}f\) in \(\mathcal{C}(\mathrm{U};\mathrm{F})\) for all \(\alpha\in A\). Let m be an integer such that \(0\leqslant m\leqslant r\) . By induction on m, it follows from th. 1 of FVR, II, p. 2, that the map \(g_{0}\) is of class \(C^{m}\) and that \(g_{\alpha}=\partial^{\alpha}g_{0}\) for all \(\alpha\in\mathbf{N}^{n}\) with \(|\alpha|\leqslant m\) . Thus \(g_{0}\) belongs to the space \(\mathcal{C}^{r}(\mathrm{U};\mathrm{F})\) and its image under u is \((g_{\alpha})_{\alpha\in\mathrm{A}}\) . This proves that the image of u is closed, and thus establishes assertion a).

The assertion \(a\) ) implies that the space \(\mathcal{C}^{r}(\mathrm{U};\mathrm{F})\) is isomorphic to its image in \(\mathcal{C}(\mathrm{U};\mathrm{F})^{\mathrm{A}}\) ; since this image is closed and the space \(\mathcal{C}(\mathrm{U};\mathrm{F})^{\mathrm{A}}\) is complete (TG, X, p. 9, cor. 3 of th. 2 and TG, II, p. 17, prop. 10), the space \(\mathcal{C}^{r}(\mathrm{U};\mathrm{F})\) is complete.

PROPOSITION 8. --- Suppose that F is finite-dimensional. Let s be an integer such that \(0 \leqslant s < r\) and let V be a relatively compact open subset of U. The linear map \(f \mapsto f|V\) from \(\mathcal{C}^{r}(U;F)\) into \(\mathcal{C}^{s}(V;F)\) is compact.

Let W be the set of functions \(f \in \mathcal{C}^{r}(\mathrm{U};\mathrm{F})\) whose partial derivatives of order \(\leqslant r\) take at every point of \(\overline{V}\) a value of norm less than 1. The set W is a neighborhood of 0 in the space \(\mathcal{C}^{r}(\mathrm{U};\mathrm{F})\) . Let \(\alpha\) be a multi-index such that \(|\alpha| \leqslant s\) . Consider the set H of functions of the form \((\partial^{\alpha}f)|\mathrm{V}\) for f in W. The set H is an equicontinuous subset of \(\mathcal{C}(\mathrm{V};\mathrm{F})\) by the mean value theorem (VAR, R, 2.2.3). Moreover, for every \(x \in V\) , the image of H under the map \(g \mapsto g(x)\) is a bounded subset, hence relatively compact, of F. By the Ascoli theorem (TG, X, p. 18, cor. 2), the set H is relatively compact in \(\mathcal{C}(\mathrm{V};\mathrm{F})\) . This proves that the linear map \(f \mapsto (\partial^{\alpha}f)|\mathrm{V}\) from \(\mathcal{C}^{r}(\mathrm{U};\mathrm{F})\) into \(\mathcal{C}(\mathrm{V};\mathrm{F})\) is compact. The proposition then follows from Lemma 3 and Remarks 5 and 6 of III, p. 3.

\subsection{Restriction of differentiable sections of a vector bundle}

Let r be a positive integer, X a differentiable manifold of class \(C^{r}\) locally of finite dimension, and E a vector bundle (real or complex),

with base X and of class \(C^{r}\). For every open set U of X, denote by \(\mathcal{S}^{r}(\mathrm{U};\mathrm{E})\) the vector space (denoted \(\mathcal{S}_{\mathrm{E}}^{r}(\mathrm{U})\) in VAR, R1, 7.4, p. 74) of \(C^{r}\) sections of E over U, equipped with the topology of \(C^{r}\) compact convergence.

Lemma 4. --- Let \(\mathcal{U}\) be an open cover of X. The map \(u: f \mapsto (f|U)_{U \in \mathcal{U}}\) from \(\mathcal{S}^{r}(X;E)\) into \(\prod_{U \in \mathcal{U}} \mathcal{S}^{r}(U;E)\) is linear, injective, continuous, and strict. Its image is closed.

The map \(u\) is linear, injective, and continuous. By virtue of TG, IX, p. 43, prop. 1 and p. 48, cor. 1, every compact subset of X has a finite cover \((\mathrm{C}_i)_{i\in \mathrm{I}}\) such that, for each \(i\in \mathrm{I}\) , the set \(\mathrm{C}_i\) is a compact subset of one of the open sets of the covering \(\mathcal{U}\) . It follows that the linear map \(u\) is strict. Finally, its image is closed, since it consists of the families \((f_{\mathrm{U}})_{\mathrm{U}\in \mathcal{U}}\) such that \(f_{\mathrm{U}}\) and \(f_{\mathrm{V}}\) coincide in \(\mathrm{U}\cap \mathrm{V}\) for all U and V in \(\mathcal{U}\) .

PROPOSITION 9. --- The space \(\mathcal{S}^r(\mathrm{X};\mathrm{E})\) is complete.

Suppose first that there exists an integer \(n \geqslant 0\) and a Banach space F such that X is an open subset of \(R^{n}\) and E is the trivial vector bundle \(X \times F\) with base X and fiber F. In this case, the topological vector space \(\mathcal{S}^{r}(X;E)\) is isomorphic to \(\mathcal{C}^{r}(X;F)\) , and the result follows from lemma 3 of III, p. 34.

In the general case, let \(\mathcal{U}\) be an open covering of X by coordinate domains such that for every \(\mathrm{U} \in \mathcal{U}\) , the restriction of E to U is trivializable. By the lemma above, the space \(\mathcal{S}^r(\mathrm{X};\mathrm{E})\) is isomorphic to the image of the linear map \(f \mapsto (f|\mathrm{U})_{\mathrm{U} \in \mathcal{U}}\) , which is closed in the product of the spaces \(\mathcal{S}^r(\mathrm{U};\mathrm{E})\) for U in \(\mathcal{U}\) . By the preceding case, each of the spaces \(\mathcal{S}^r(\mathrm{U};\mathrm{E})\) is complete, hence their product is complete (TG, II, p. 17, prop. 10). Consequently, \(\mathcal{S}^r(\mathrm{X};\mathrm{E})\) is complete.

Remark. --- Let \(\mathcal{U}\) be an open covering of X by coordinate domains \(c_{\mathrm{U}} = (\mathrm{U}, \varphi_{\mathrm{U}}, \mathbf{R}^{n_{\mathrm{U}}})\) such that for every \(U \in \mathcal{U}\) , the restriction of E to U is trivializable of type \(F_{U}\) , where \(F_{U}\) is a Banach space. For every \(U \in \mathcal{U}\) , one identifies a section f of E over U with a map \(f_{U}\) from \(\varphi_{\mathrm{U}}(\mathrm{U})\) into \(F_{U}\) . It follows from the proof of prop. 9 that the topology of the space \(\mathcal{S}^{r}(\mathrm{X};\mathrm{E})\) is defined by the family of seminorms \(p_{U,C,\alpha}\) such that

\[
p _ {\mathrm{U,C}, \alpha} (f) = \sup _ {x \in \mathrm{C}} \| (\partial^ {\alpha} f _ {\mathrm{U}}) (x) \|,
\]

where U ranges over \(\mathcal{U}\) , C ranges over the set of compact subsets of \(\varphi_{\mathrm{U}}(\mathrm{U})\) and \(\alpha \in \mathbf{N}^{n_{\mathrm{U}}}\) ranges over the set of multi-indices such that \(|\alpha| \leqslant r\) .

PROPOSITION 10. --- Suppose that the vector bundle E is locally of finite rank. Let s be an integer such that \(0 \leqslant s < r\) and Y a relatively compact open subset of X. The linear map \(f \mapsto f|Y\) from \(\mathcal{S}^{r}(X;E)\) into \(\mathcal{S}^{s}(Y;E)\) is compact.

Suppose first that there exists an integer \(n \geqslant 0\) and a finite-dimensional Banach space F such that X is an open subset of \(R^{n}\) and E is the trivial vector bundle \(X \times F\) with base X and fiber F. In this case, the topological vector spaces \(\mathcal{S}^{r}(X;E)\) and \(\mathcal{S}^{s}(Y;E)\) are isomorphic to \(\mathcal{C}^{r}(X;F)\) and \(\mathcal{C}^{s}(Y;F)\) , respectively, and prop. 10 is a consequence of prop. 8.

Let us pass to the general case. Let C be a compact subset of X containing Y. For every point x of C, choose an open neighborhood \(U(x)\) of x in X which is a coordinate domain over which the vector bundle E is trivializable. Choose moreover a relatively compact open neighborhood \(V(x)\) of x in \(U(x)\) . Since the set C is compact, it is covered by a finite family \((\mathrm{V}(x_{1}),\ldots,\mathrm{V}(x_{m}))\) of such open sets. For every i, set \(U_{i}=U(x_{i})\) and \(Y_{i}=V(x_{i})\cap Y\) . Then \(Y_{i}\) is a relatively compact open subset of \(U_{i}\) and we have \(Y=Y_{1}\cup\cdots\cup Y_{m}\) . The linear maps \(f\mapsto f|U_{i}\) from \(\mathcal{S}^{r}(X;E)\) into \(\mathcal{S}^{r}(U_{i};E)\) are continuous and the linear maps \(g\mapsto g|Y_{i}\) from \(\mathcal{S}^{r}(U_{i};E)\) into \(\mathcal{S}^{s}(Y_{i};E)\) are compact by the first part of the proof. The linear maps \(f\mapsto f|Y_{i}\) from \(\mathcal{S}^{r}(X;E)\) into \(\mathcal{S}^{s}(Y_{i};E)\) are therefore compact (prop. 3 of III, p. 5). Taking into account Lemma 4, applied to the covering of Y by the open sets \(Y_{i}\) , and Remarks 5 and 6 of III, p. 3, the linear map \(f\mapsto f|Y\) from \(\mathcal{S}^{r}(X;E)\) into \(\mathcal{S}^{s}(Y;E)\) is compact, which completes the proof.

\subsection{Restriction of analytic sections of a vector bundle}

Let X be a complex analytic manifold locally of finite dimension and E an analytic vector bundle with base X (VAR, p. 35 and VAR, p. 70). Denote by \(\mathcal{S}^{\omega}(\mathrm{X};\mathrm{E})\) the vector space of analytic sections of E over X, equipped with the topology of compact convergence.

Lemma 5. --- Let \(X_{0}\) be the underlying real analytic manifold of X and \(E_{0}\) the complex vector bundle over \(X_{0}\) underlying E.

\begin{enumerate}
\def\labelenumi{\alph{enumi})}
\item
  The vector subspace \(\mathcal{S}^{\omega}(\mathrm{X};\mathrm{E})\) of \(\mathcal{S}^{0}(\mathrm{X}_{0};\mathrm{E}_{0})\) is closed, and the injection of \(\mathcal{S}^{\omega}(\mathrm{X};\mathrm{E})\) into \(\mathcal{S}^{0}(\mathrm{X}_{0};\mathrm{E}_{0})\) is continuous and strict;
\item
  The canonical injection of \(\mathcal{S}^{\omega}(\mathrm{X};\mathrm{E})\) into \(\mathcal{S}^{1}(\mathrm{X}_{0};\mathrm{E}_{0})\) is continuous.
\end{enumerate}

Suppose first that there exists an integer \(n \geqslant 0\) and a Banach space F such that X is an open subset of \(C^{n}\) and E is the trivial vector bundle \(X \times F\) . In this case, the topological vector spaces \(\mathcal{S}^{\omega}(X;E)\) , \(\mathcal{S}^{0}(X_{0};E_{0})\) and \(\mathcal{S}^{1}(X_{0};E_{0})\) are isomorphic to \(\mathcal{C}^{\omega}(X;F)\) , \(\mathcal{C}^{0}(X_{0};F)\) and \(\mathcal{C}^{1}(X_{0};F)\) respectively, and the space \(\mathcal{S}^{0}(X_{0};E_{0})\) is metrizable (TG, X, p. 20, cor. of prop. 1). The lemma then follows from VAR, 3.3.2, p. 28.

Let \(\mathfrak{F}\) be a filter on \(\mathcal{S}^{\omega}(\mathrm{X};\mathrm{E})\) which converges in the space \(\mathcal{S}^{0}(\mathrm{X}_{0};\mathrm{E}_{0})\) toward a limit \(f\). We must show that \(f\) belongs to \(\mathcal{S}^{\omega}(\mathrm{X};\mathrm{E})\) and that the filter \(\mathfrak{F}\) converges to \(f\) in the space \(\mathcal{S}^{1}(\mathrm{X}_{0};\mathrm{E}_{0})\). This statement is local in nature, and therefore follows from the first part of the proof.

PROPOSITION 11. --- Suppose that the vector bundle E is locally of finite rank. Let Y be a relatively compact open subset of X. The restriction map \(f \mapsto f|Y\) from \(\mathcal{S}^{\omega}(X;E)\) into \(\mathcal{S}^{\omega}(Y;E)\) is compact.

With the notations of Lemma 5, we have a commutative diagram

\[
\begin{array}{c c c} \mathcal {S} ^ {\omega} (\mathrm{X}; \mathrm{E}) & \stackrel {{i}} {{\longrightarrow}} & \mathcal {S} ^ {1} (\mathrm{X} _ {0}; \mathrm{E} _ {0}) \\ \Big \downarrow^ {u} & & \Big \downarrow^ {v} \\ \mathcal {S} ^ {\omega} (\mathrm{Y}; \mathrm{E}) & \stackrel {{j}} {{\longrightarrow}} & \mathcal {S} ^ {0} (\mathrm{Y} _ {0}; \mathrm{E} _ {0}) \end{array}\tag{3}
\]

where i and j are the canonical injections and u, v the restriction maps. The map i is continuous (Lemma 5, b)), and the map v is compact (Prop. 10). Since the canonical injection j is continuous, strict, and has closed image (Lemma 5, a)), the map u is compact (Remark 5 of III, p. 3).

COROLLARY. --- Let X be a compact complex analytic manifold and E an analytic vector bundle over X, locally of finite rank. The vector space \(\mathcal{S}^{\omega}(X;E)\) is of finite dimension.

Being compact, the analytic manifold X is locally of finite dimension. By prop. 11, the identity map of \(\mathcal{S}^{\omega}(\mathrm{X};\mathrm{E})\) is a compact linear map. This implies that the space \(\mathcal{S}^{\omega}(\mathrm{X};\mathrm{E})\) is of finite dimension (Remark 3 of III, p. 2).

\section{FREDHOLM ENDOMORPHISMS AND RIESZ ENDOMORPHISMS}

\subsection{Strict morphisms and finite-rank linear maps}

PROPOSITION 1. --- Let E and F be locally convex spaces, \(E_{1}\) a closed subspace of finite codimension of E and u a continuous linear map from E into F. For u to be a strict morphism with closed image, it is necessary and sufficient that \(u|E_{1}\) be one.

Suppose first that the linear map \(u|E_1\) is injective. We then have \(E_1 \cap \text{Ker}(u) = \{0\}\) so that \(\text{Ker}(u)\) is finite-dimensional. Let S be a vector subspace of E supplementary to \(E_1 + \text{Ker}(u)\) . The space E is the topological direct sum of \(E_1\) , \(\text{Ker}(u)\) and S (EVT, I, p. 15, cor. 4 and prop. 3). If \(u\) is a strict morphism with closed image, it defines by restriction an isomorphism from \(E_1 \oplus S\) onto a closed vector subspace of F, and a fortiori an isomorphism from \(E_1\) onto a closed vector subspace of F. Conversely, suppose that \(u\) defines by restriction an isomorphism from \(E_1\) onto \(u(E_1)\) and that \(u(E_1)\) is closed in F. We have \(u(E) = u(E_1) \oplus u(S)\) . It follows that \(u(E)\) is the topological direct sum of \(u(E_1)\) and \(u(S)\) , and is closed in F (loc. cit.). Since \(u(E_1)\) is closed and since we have \(u(E_1) \cap u(S) = \{0\}\) , the space \(u(S)\) is separated and \(u\) defines by restriction an isomorphism from S onto \(u(S)\) (EVT, I, p. 15, cor. of prop. 3), hence also of \(E_1 \oplus S\) onto \(u(E)\) . This proves that \(u\) is a strict morphism with closed image.

Moving to the general case, set \(N = E_{1} \cap \operatorname{Ker}(u)\) and G = E/N. The locally convex space \(E_{1}/N\) is identified with a closed vector subspace of finite codimension \(G_{1}\) of G. Let us denote by \(v\colon G \to F\) the continuous linear map deduced from u by passing to the quotient. For \(u\) (resp. \(u|E_{1}\)) to be a strict morphism with closed image, it is necessary and sufficient that \(v\) (resp. \(v|G_{1}\)) be one. This reduces us to the case already treated.

COROLLARY 1. --- Suppose F is separated. Let \(v \in \mathcal{L}^{\mathrm{f}}(\mathrm{E};\mathrm{F})\) . If u is a strict morphism with closed image from E into F, then the same holds for \(u + v\) .

Since F is separated, the kernel of v is a closed vector subspace of E; it has finite codimension in E and \(u + v\) has the same restriction as \(u\) on \(\operatorname{Ker}(v)\) . The corollary thus follows from the proposition.

COROLLARY 2. --- Let T be a separated finite-dimensional vector subspace of F. Denote by \(\pi\colon\mathrm{F}\to\mathrm{F}/\mathrm{T}\) the canonical surjection. For u to be a strict morphism with closed image, it is necessary and sufficient that \(\pi\circ u\) be one.

The identity map of T onto itself extends to a continuous map \(q \colon \mathrm{F} \to \mathrm{T}\) (EVT, II, p. 26, remark). The kernel S of \(q\) is a closed topological complement of T. Let \(p \colon \mathrm{F} \to \mathrm{F}\) be the projector with image S associated with the decomposition \(\mathrm{F} = \mathrm{T} \oplus \mathrm{S}\) . For \(\pi \circ u\) to be a strict morphism with closed image, it is necessary and sufficient that \(p \circ u\) be one. Moreover, \(p \circ u\) and \(u\) have the same restriction to \(u^{-1}(\mathrm{S})\) , which is a closed vector subspace of finite codimension in E, and the assertion follows from the proposition.

\subsection{Fredholm Applications}

Let E and F be locally convex spaces. In this section, we denote by \(u \equiv v\) the congruence modulo \(\mathcal{L}^{\mathrm{f}}(\mathrm{E};\mathrm{F})\) of the elements u and v of \(\mathcal{L}(\mathrm{E};\mathrm{F})\) .

If G is a locally convex space, and if \(u'\) and \(v'\) are elements of \(\mathcal{L}(\mathrm{F};\mathrm{G})\) , the relations \(u \equiv v\) and \(u' \equiv v'\) imply the relation \(u' \circ u \equiv v' \circ v\) .

We say that an element \(w\) of \(\mathcal{L}(\mathrm{F};\mathrm{E})\) is a quasi-inverse of the element \(u\) of \(\mathcal{L}(\mathrm{E};\mathrm{F})\) if we have \(w \circ u \equiv 1_{\mathrm{E}}\) and \(u \circ w \equiv 1_{\mathrm{F}}\) .

Suppose that w is a quasi-inverse of u. If \(u_{1}\) is an element of \(\mathcal{L}(\mathrm{E};\mathrm{F})\) and \(w_{1}\) is an element of \(\mathcal{L}(\mathrm{F};\mathrm{E})\) such that \(u_{1} \equiv u\) and \(w_{1} \equiv w\) , then \(w_{1}\) is a quasi-inverse of \(u_{1}\) since \(w \circ u \equiv w_{1} \circ u_{1}\) and \(u \circ w \equiv u_{1} \circ w_{1}\) .

If w and \(w_{1}\) are quasi-inverses of u, then \(w_{1} \equiv w\) since

\[
w _ {1} = 1 _ {\mathrm{E}} \circ w _ {1} \equiv w \circ u \circ w _ {1} \equiv w \circ 1 _ {\mathrm{F}} = w.
\]

DEFINITION 1. --- Let E and F be locally convex spaces. We say that an element u of \(\mathcal{L}(\mathrm{E};\mathrm{F})\) is a Fredholm map \(^{(1)}\) if it possesses a quasi-inverse. A Fredholm map from E to E is called a Fredholm endomorphism of E.

We shall denote by \(\mathcal{F}(\mathrm{E};\mathrm{F})\) the set of Fredholm maps from E to F, and \(\mathcal{F}(\mathrm{E})\) the set of Fredholm endomorphisms of E.

Remarks. --- Let E, F, and G be locally convex spaces, \(u: E \to F\) and \(v: F \to G\) continuous linear maps.

\begin{enumerate}
\def\labelenumi{\arabic{enumi})}
\item
  Suppose that u is a Fredholm map and denote by \(u_{1}\) a quasi-inverse of u. Since u is a quasi-inverse of \(u_{1}\) , the map \(u_{1}\) is a Fredholm map.
\item
  Suppose that u and v are Fredholm maps, and let \(u_{1}\) (resp. \(v_{1}\) ) be a quasi-inverse of u (resp. v). Then \(v \circ u\) is a Fredholm map from E to G and \(u_{1} \circ v_{1}\) is a quasi-inverse of \(v \circ u\) . Indeed, one computes
\end{enumerate}

\[
\left(u _ {1} \circ v _ {1}\right) \circ (v \circ u) = u _ {1} \circ \left(v _ {1} \circ v\right) \circ u \equiv u _ {1} \circ 1 _ {\mathrm{F}} \circ u = u _ {1} \circ u \equiv 1 _ {\mathrm{E}},
\]

\[
(v \circ u) \circ (u _ {1} \circ v _ {1}) = v \circ (u \circ u _ {1}) \circ v _ {1} \equiv v \circ 1 _ {\mathrm{F}} \circ v _ {1} = v \circ v _ {1} \equiv 1 _ {\mathrm{G}}.
\]

\begin{enumerate}
\def\labelenumi{\arabic{enumi})}
\setcounter{enumi}{2}
\tightlist
\item
  Suppose that u and \(v \circ u\) are Fredholm maps, and let \(w_{1}\) be a quasi-inverse of \(v \circ u\) . Then v is a Fredholm map and \(u \circ w_{1}\) is a quasi-inverse of v.
\end{enumerate}

Indeed, \(w_{1}\) is a Fredholm operator by the first remark. Let \(u_{1}\) be a quasi-inverse of u; by the second remark, \(u \circ w_{1}\) is a Fredholm map and \((v \circ u) \circ u_{1}\) is a quasi-inverse. We have \(u \circ u_{1} \equiv 1_{F}\) , whence \(v \circ u \circ u_{1} \equiv v\) ; this proves the assertion.

\begin{enumerate}
\def\labelenumi{\arabic{enumi})}
\setcounter{enumi}{3}
\tightlist
\item
  Suppose that v and \(v \circ u\) are Fredholm maps, and let \(w_{1}\) be a quasi-inverse of \(v \circ u\) . Then u is a Fredholm operator and \(w_{1} \circ v\) is a quasi-inverse of u.
\end{enumerate}

The proof is analogous to that of the preceding remark.

Lemma 1. --- Let E and F be locally convex spaces and u a Fredholm map from E into F. The kernel and cokernel of u are finite-dimensional.

Let \(v \colon \mathrm{F} \to \mathrm{E}\) be a quasi-inverse of \(u\) . The kernel of \(u\) is contained in the image of the finite-rank linear map \(1_{\mathrm{E}} - v \circ u\) , hence is finite-dimensional. The image of \(u\) contains the kernel of the finite-rank linear map \(1_{\mathrm{F}} - u \circ v\) , hence is of finite codimension in F.

PROPOSITION 2. — Let E and F be separated locally convex spaces and let u be an element of \(\mathcal{L}(\mathrm{E};\mathrm{F})\) . The following properties are equivalent:

\begin{enumerate}
\def\labelenumi{(\roman{enumi})}
\item
  The map u is a Fredholm operator;
\item
  The map \(u\) is a strict morphism, its kernel is finite-dimensional, its image is closed and of finite codimension in F;
\item
  There exist closed vector subspaces \(E_{1}\) of E and \(F_{1}\) of F, both of finite codimension, such that u induces, by passing to subspaces, an isomorphism from \(E_{1}\) onto \(F_{1}\) ;
\item
   There exist decompositions into topological direct sums \(E = E_{1} \oplus E_{2}\) and \(F = F_{1} \oplus F_{2}\) such that \(E_{2}\) and \(F_{2}\) are finite-dimensional, that u vanishes on \(E_{2}\) and defines by passing to subspaces an isomorphism of \(E_{1}\) onto \(F_{1}\) .
\item
  \(\Longrightarrow\) (iii): Let v be a quasi-inverse of u, \(E_{1}\) the kernel of \(1_{E}-v\circ u\) and \(F_{1}\) that of \(1_{F}-u\circ v\) . Since the linear maps \(1_{E}-v\circ u\) and \(1_{F}-u\circ v\) are continuous and of finite rank, \(E_{1}\) and \(F_{1}\) are closed vector subspaces of finite codimension of E and F, respectively. Let \(x\in E_{1}\) . We have
\end{enumerate}

\[
(1 _ {\mathrm{F}} - u \circ v) (u (x)) = u ((1 _ {\mathrm{E}} - v \circ u) (x)) = u (0) = 0,
\]

whence \(u(x) \in \mathrm{F}_1\) . We therefore have \(u(\mathrm{E}_1) \subset \mathrm{F}_1\) ; we likewise have \(v(\mathrm{F}_1) \subset \mathrm{E}_1\) . The continuous linear maps \(u_1: \mathrm{E}_1 \to \mathrm{F}_1\) and \(v_1: \mathrm{F}_1 \to \mathrm{E}_1\) deduced from \(u\) and \(v\) are then mutually inverse isomorphisms, since \(v \circ u\) and \(1_{\mathrm{E}}\) (resp. \(u \circ v\) and \(1_{\mathrm{F}}\) ) coincide on \(\mathrm{E}_1\) (resp. on \(\mathrm{F}_1\) ).

\begin{enumerate}
\def\labelenumi{(\roman{enumi})}
\setcounter{enumi}{2}
\item
  \(\Longrightarrow\) (ii): Let \(E_{1}\) and \(F_{1}\) satisfy the hypothesis of (iii). We have \(\mathrm{E}_{1}\cap\mathrm{Ker}(u)=\{0\}\) and \(\mathrm{F}_{1}\subset\mathrm{Im}(u)\) , hence \(\mathrm{Ker}(u)\) is finite-dimensional and \(\mathrm{Im}(u)\) is closed and of finite codimension in F. It follows from prop. 1 of III, p. 39 that the map u is a strict morphism.
\item
   \(\Longrightarrow\) (iv): Suppose condition (ii) is satisfied. The closed vector subspace \(E_{2} = \operatorname{Ker}(u)\) of E is finite-dimensional, and there exists a vector subspace \(E_{1}\) of E topologically complementary to \(E_{2}\) (EVT, II, p. 27, cor. 2). The vector subspace \(F_{1} = \operatorname{Im}(u)\) is closed and of finite codimension, and it admits a topological complement \(F_{2}\) in F. By prop. 1 of III, p. 39, the map \(u|E_{1}\) is a strict morphism, hence u induces an isomorphism of \(E_{1}\) onto \(F_{1}\) .
\item
   \(\Longrightarrow\) (i): Under the hypotheses of (iv), the linear map from F to E that coincides with \(u_1^{-1}\) on \(\mathrm{F}_1\) and is zero on \(\mathrm{F}_2\) is a quasi-inverse of \(u\) .
\end{enumerate}

Remark 5. --- Let E, F be separated locally convex spaces and \(u\colon\mathrm{E}\to\mathrm{F}\) a Fredholm operator. If \(u\) is bijective, then \(u\) is an isomorphism (indeed, \(u\) is a strict morphism by Prop. 2, (ii)).

\subsection{Index of a Fredholm map}

Let E, F, and G be locally convex spaces.

Let \(u \colon \mathrm{E} \to \mathrm{F}\) be a Fredholm operator. The vector spaces \(\operatorname{Ker}(u)\) and \(\operatorname{Coker}(u)\) are of finite dimension (Lemma 1 of III, p. 41). Recall that the integer

\[
\dim \operatorname{Coker} (u) - \dim \operatorname{Ker} (u) = \operatorname{codim} _ {\mathrm{F}} \operatorname{Im} (u) - \dim \operatorname{Ker} (u) \tag {1}
\]

is called the index of u and is denoted ind(u) (A, V, p. 126).

If \(u \colon \mathrm{E} \to \mathrm{F}\) and \(v \colon \mathrm{F} \to \mathrm{G}\) are Fredholm maps, then so is \(v \circ u\) (III, p. 41, remark 2), and one has (A, V, p. 127, lemma 2)

\[
\operatorname{ind} (v \circ u) = \operatorname{ind} (v) + \operatorname{ind} (u).\tag{2}
\]

Suppose E and F are separated, and let \(u\colon\mathrm{E}\to\mathrm{F}\) be a Fredholm operator; let us adopt the notations of condition (iv) of Prop. 2 of III, p. 42. We then have \(\mathrm{ind}(u)=\mathrm{dim}(\mathrm{F}_{2})-\mathrm{dim}(\mathrm{E}_{2})\) . Equip the dual of each of these spaces with the weak topology (resp. the compact-open topology, the topology of bounded convergence). Then \(\mathrm{E}^{\prime}\) is identified with the topological direct sum of \(\mathrm{E}_{1}^{\prime}\) and \(\mathrm{E}_{2}^{\prime}\) , and \(\mathrm{F}^{\prime}\) is identified with that of \(\mathrm{F}_{1}^{\prime}\) and \(\mathrm{F}_{2}^{\prime}\) , and \(^{t}u\) induces an isomorphism of \(\mathrm{F}_{1}^{\prime}\) onto \(\mathrm{E}_{1}^{\prime}\) and vanishes on \(\mathrm{F}_{2}^{\prime}\) . Thus the transpose \(^{t}u\colon\mathrm{F}^{\prime}\to\mathrm{E}^{\prime}\) is a Fredholm map (loc. cit.). The kernel of \(^{t}u\) is \(\mathrm{F}_{2}^{\prime}\) , and its dimension is that of \(\mathrm{F}_{2}\) , that is, that of the cokernel of \(u\) . Hence

\[
\operatorname{ind} (u) = \dim \operatorname{Ker} (^ {t} u) - \dim \operatorname{Ker} (u).\tag{3}
\]

Moreover, the image of \(^{t}u\) is \(E_{1}^{\prime}\) , and the dimension of the cokernel of \(^{t}u\) is therefore equal to the dimension of \(E_{2}^{\prime}\) , which is that of the kernel \(E_{2}\) of u. We deduce from this

\[
\operatorname{ind} (^ {t} u) = - \operatorname{ind} (u).\tag{4}
\]

Suppose E and F are separated, and let \(u\colon\mathrm{E}\to\mathrm{F}\) be a Fredholm operator of index 0. Then \(u\) is a strict morphism by prop. 2 of III, p. 42. Since \(\dim\operatorname{Ker}(u)=\operatorname{codim}_{\mathrm{F}}\operatorname{Im}(u)\) , the map \(u\) is an isomorphism as soon as it is injective or surjective.

PROPOSITION 3. --- Let E and F be locally convex spaces and \(u \in \mathcal{L}(\mathrm{E};\mathrm{F})\) . Let \(\mathrm{E}_{1}\) (resp. \(\mathrm{F}_{1}\) ) be a closed subspace of finite codimension of E (resp. F). Suppose that u maps \(\mathrm{E}_{1}\) into \(\mathrm{F}_{1}\) and denote by \(u_{1} \in \mathcal{L}(\mathrm{E}_{1};\mathrm{F}_{1})\) the map which coincides with u on \(\mathrm{E}_{1}\) . For u to be a Fredholm operator, it is necessary and sufficient that \(u_{1}\) be one. We then have

\[
\operatorname{ind} (u) - \operatorname{ind} \left(u _ {1}\right) = \operatorname{codim} _ {\mathrm{F}} \left(\mathrm{F} _ {1}\right) - \operatorname{codim} _ {\mathrm{E}} \left(\mathrm{E} _ {1}\right).\tag{5}
\]

Let \(i\colon \mathrm{E}_{1}\to \mathrm{E}\) and \(j\colon \mathrm{F}_{1}\to \mathrm{F}\) be the canonical injections. These are Fredholm maps, and one has

\[
\operatorname{ind} (i) = \operatorname{codim} _ {\mathrm{E}} (\mathrm{E} _ {1}), \quad \operatorname{ind} (j) = \operatorname{codim} _ {\mathrm{F}} (\mathrm{F} _ {1}).
\]

Since \(j \circ u_{1} = u \circ i\) , one sees that u is a Fredholm map if and only if \(u_{1}\) is one (remarks 3 and 4 of III, p. 41). If this is so, we have

\[
\operatorname{ind} (j) + \operatorname{ind} (u _ {1}) = \operatorname{ind} (j \circ u _ {1}) = \operatorname{ind} (u \circ i) = \operatorname{ind} (u) + \operatorname{ind} (i),
\]

whence formula (5).

PROPOSITION 4. --- Let E and F be separated locally convex spaces, \(u: \mathrm{E} \to \mathrm{F}\) a Fredholm operator, and \(\widehat{u}: \widehat{\mathrm{E}} \to \widehat{\mathrm{F}}\) the extension of u to the completions. Then \(\widehat{u}\) is a Fredholm operator, and one has \(\operatorname{Ker}(\widehat{u}) = \operatorname{Ker}(u)\) and \(\operatorname{ind}(\widehat{u}) = \operatorname{ind}(u)\) .

Let us adopt the notation of condition (iv) of Prop. 2 of III, p. 42. Since the vector spaces \(E_{2}\) and \(F_{2}\) are finite-dimensional, they are complete. The completion of \(E_{1}\) (resp. \(F_{1}\) ) is identified with the closure of \(E_{1}\) in \(\widehat{E}\) (resp. of \(F_{1}\) in \(\widehat{F}\) ), and \(\widehat{E}\) (resp. \(\widehat{F}\) ) is the topological direct sum of \(\widehat{E}_{1}\) and \(E_{2}\) (resp. \(\widehat{F}_{1}\) and \(F_{2}\) ). The linear map \(\widehat{u}\) defines by restriction an isomorphism of \(\widehat{E}_{1}\) onto \(\widehat{F}_{1}\) and vanishes on \(E_{2}\) . The proposition then follows from implication (iv)⇒(i) of loc. cit.

\subsection{Riesz endomorphisms}

Let E be a vector space and let u be an endomorphism of E.

The sequence \((\operatorname{Ker}(u^{n}))_{n\in\mathbf{N}}\) of vector subspaces of E is increasing; their union is a subspace of E stable under u, called the nilspace of u.

The sequence \((\mathrm{Im}(u^{n}))_{n\in\mathbf{N}}\) of vector subspaces of E is decreasing; their intersection is a subspace of E stable under u, called the conilspace of u.

Lemma 2. --- Let E be a vector space and let u be an endomorphism of E which has an index (A, V, p. 126).

If two of the following properties are satisfied, then so is the third:

\begin{enumerate}
\def\labelenumi{(\roman{enumi})}
\item
  The endomorphism u has index 0;
\item
  The nilespace N of u is finite-dimensional;
\item
  The conilespace I of u has finite codimension.
\end{enumerate}

For every integer \(n \geqslant 0\) , the endomorphism \(u^{n}\) has an index, equal to \(n \operatorname{ind}(u)\) (A, V, p. 127, lemma 2). The sequence \((\dim(\operatorname{Ker}(u^{n})))_{n \in \mathbf{N}}\) of natural integers is increasing; for it to be stationary, it is necessary and sufficient that the nilespace of u be finite-dimensional. The sequence \((\operatorname{codim}_{\operatorname{E}}(\operatorname{Im}(u^{n})))_{n \in \mathbf{N}}\) of natural integers is increasing; for it to be stationary, it is necessary and sufficient that the conilespace of u be finite-dimensional. Suppose condition (i) is satisfied; the relation

\[
\dim (\operatorname{Ker} (u ^ {n})) - \operatorname{codim} _ {\mathrm{E}} (\operatorname{Im} (u ^ {n})) = \operatorname{ind} (u ^ {n}) = n \operatorname{ind} (u) = 0
\]

then implies the equivalence of conditions (ii) and (iii). Conversely, if conditions (ii) and (iii) are satisfied, this formula implies that the sequence \((n\operatorname{ind}(u))_{n\in\mathbf{N}}\) is stationary, so u has index zero.

Lemma 3. --- Let E be a vector space and let u be an endomorphism of E which has an index. Suppose that the index of u is zero. Let N be the nilspace of u and I its conilspace.

\begin{enumerate}
\def\labelenumi{\alph{enumi})}
\item
  The space E is then the direct sum of the subspaces N and I (Weyr–Fitting decomposition, cf. A, VIII, §2, no. 2, p. 25).
\item
  The endomorphism u induces, by passing to subspaces, a nilpotent endomorphism of N and an automorphism \(u_{I}\) of I.
\item
  Let v denote the endomorphism of E which is zero on N and which coincides with \(u_{I}^{-1}\) on I. We have \(u \circ v = v \circ u\) and the endomorphism \(1_{E} - u \circ v\) is the projector with image N and kernel I.
\item
  Every endomorphism of E that commutes with u stabilizes N and I, and commutes with v.
\end{enumerate}

The endomorphism \(u\) satisfies the properties of the preceding lemma.

Let \(n \in \mathbf{N}\) such that \(\operatorname{Ker}(u^{n}) = N\) . Then \(\operatorname{Ker}(u^{m}) = \operatorname{Ker}(u^{n})\) for every integer \(m \geqslant n\) ; since \(\operatorname{ind}(u^{m}) = m \operatorname{ind}(u) = 0\) , this entails \(\operatorname{Im}(u^{m}) = \operatorname{Im}(u^{n})\) for \(m \geqslant n\) , hence \(\operatorname{I} = \operatorname{Im}(u^{n})\) . It follows from A, VIII, p. 25, prop. 2 that assertions a) and b) are satisfied.

For every \(x \in N\) , we have \(v(x) = 0\) and \(u(x) \in N\) , hence \(u(v(x)) = 0\) and \(v(u(x)) = 0\) . For every \(x \in I\) , we have \(u(v(x)) = v(u(x)) = x\) . This implies that \(u \circ v = v \circ u\) and that \(1_{E} - u \circ v\) is the projector with image N and kernel I.

Let w be an endomorphism of E that commutes with u. Let \(n \in \mathbf{N}\) such that \(N = \operatorname{Ker}(u^{n})\) and \(I = \operatorname{Im}(u^{n})\) . Let \(x \in N\) ; we have \(u^{n}(w(x)) = w(u^{n}(x)) = 0\) , hence \(w(x) \in N\) . Let \(x \in I\) ; there exists \(y \in E\) such that \(x = u^{n}(y)\) ; then, \(w(x) = w(u^{n}(y)) = u^{n}(w(y)) \in I\) . This proves that w stabilizes N and I.

Let us finally prove that \(w\) and \(v\) are permutable. For every \(x\) in \(N\) , we have \(v(w(x)) = 0 = w(v(x))\) . Let \(x \in I\) ; the elements \(v(w(x))\) and \(w(v(x))\) of \(E\) belong to \(I\) and their images under \(u\) are both equal to \(w(x)\) , hence they are equal. By linearity, this implies that \(v\) and \(w\) commute.

Remark. --- Suppose there exists an integer \(n \geqslant 0\) such that \(\operatorname{Ker}(u^{n}) = \operatorname{Ker}(u^{n+1})\) (resp. \(\operatorname{Im}(u^{n}) = \operatorname{Im}(u^{n+1})\) ); one then has \(\operatorname{Ker}(u^{m}) = \operatorname{Ker}(u^{n})\) (resp. \(\operatorname{Im}(u^{m}) = \operatorname{Im}(u^{n})\) ) for every integer \(m \geqslant n\) .

Suppose that \(u\) has index zero. For every integer \(n \geqslant 0\) , we have

\[
n \operatorname{ind} (u) = \operatorname{ind} \left(u ^ {n}\right) = \operatorname{codim} _ {\mathrm{E}} \left(\operatorname{Im} \left(u ^ {n}\right)\right) - \dim \left(\operatorname{Ker} \left(u ^ {n}\right)\right).
\]

The conditions \(\operatorname{Ker}(u^n) = \operatorname{Ker}(u^{n+1})\) and \(\operatorname{Im}(u^n) = \operatorname{Im}(u^{n+1})\) are therefore equivalent. When they are satisfied, the vector space E is the direct sum of \(\operatorname{Ker}(u^n)\) and of \(\operatorname{Im}(u^n)\) (A, VIII, p. 25, prop. 2), the space \(\operatorname{Ker}(u^n)\) is finite-dimensional, and \(u\) induces, by passage to subspaces, an automorphism of \(\operatorname{Im}(u^n)\) .

DEFINITION 2. --- Let E be a separated locally convex space. A Riesz endomorphism of E is any Fredholm endomorphism of E whose nilespace is finite-dimensional and whose conilespace is of finite codimension.

PROPOSITION 5. --- Let E be a separated locally convex space. Every Riesz endomorphism of E has index zero.

This follows from the definition and Lemma 2.

Examples. --- 1) Every automorphism of E is a Riesz endomorphism of E. If E is finite-dimensional, every element of \(\mathcal{L}(\mathrm{E})\) is a Riesz endomorphism of E.

\begin{enumerate}
\def\labelenumi{\arabic{enumi})}
\setcounter{enumi}{1}
\item
   If \((\mathrm{E}_{i})_{i\in\mathrm{I}}\) is a finite family of separated locally convex spaces and \(u_{i}\) is an element of \(\mathcal{L}(\mathrm{E}_{i})\) for all i in I, the endomorphism \(u=\bigoplus_{i\in\mathrm{I}}u_{i}\) of the locally convex space \(\bigoplus_{i\in\mathrm{I}}E_{i}\) is a Riesz endomorphism if and only if \(u_{i}\) is one for all i.
\item
   Let E be a separated locally convex space over R, and let \(\mathrm{E}_{(\mathbf{C})}\) be its complexification. Let u be an endomorphism of E. For u to be a Riesz endomorphism, it is necessary and sufficient that its complexification \(u_{(\mathbf{C})}\) be a Riesz endomorphism of \(\mathrm{E}_{(\mathbf{C})}\) .
\end{enumerate}

PROPOSITION 6. --- Let E be a separated locally convex space and let u be a Riesz endomorphism of E. Denote by N the nilspace of u and by I its conilspace.

\begin{enumerate}
\def\labelenumi{\alph{enumi})}
\item
  The subspaces N and I of E are closed and stable under u, and the space E is their topological direct sum.
\item
  The endomorphism u defines by restriction an automorphism \(u_{I}\) of I;
\item
   The vector space N is finite-dimensional, and u defines by restriction a nilpotent endomorphism \(u_{N}\) of N;
\item
  Let v be the element of \(\mathcal{L}(\mathrm{E})\) that is zero on N and coincides with \(u_{\mathrm{I}}^{-1}\) on I. It is a Riesz endomorphism of E and a quasi-inverse of u that commutes with u. The endomorphism \(1_{\mathrm{E}} - u \circ v\) of E is the projector with image N and kernel I. Every element of \(\mathcal{L}(\mathrm{E})\) which commutes with u stabilizes N and I, and commutes with v.
\end{enumerate}

Let n be a natural number such that \(\operatorname{Ker}(u^{n}) = \operatorname{N}\) and \(\operatorname{Im}(u^{n}) = \operatorname{I}\) . Since E is separated and u is continuous, \(\operatorname{N} = \operatorname{Ker}(u^{n})\) is closed in E. Since u is a Fredholm endomorphism, so is \(u^{n}\) , and \(\operatorname{I} = \operatorname{Im}(u^{n})\) is closed (III, p. 42, prop. 2). By lemma 2, the space E is the direct sum of N and I, and assertion a) then follows from prop. 3 of EVT, I, p. 15.

Since I is closed, of finite codimension in E, and stable under u, the map \(u_{I}\colon I \to I\) induced by u is a Fredholm endomorphism (III, p. 44, prop. 3); it is also bijective, hence it is an automorphism of I (III, p. 42, prop. 2). This proves b), and the remaining assertions follow from lemma 2.

The endomorphism \(v\) defined in assertion \(d\) ) of the proposition is a quasi-inverse of \(u\) , which is called the canonical quasi-inverse of \(u\) .

Riesz endomorphisms enjoy certain of the stability properties of Fredholm maps.

PROPOSITION 7. — Let E be a separated locally convex space and let u be an endomorphism of E.

\begin{enumerate}
\def\labelenumi{\alph{enumi})}
\item
   Let \(E_{1}\) be a closed subspace of finite codimension of E, stable under u. Let \(u_{1}\) be the element of \(\mathcal{L}(E_{1})\) which coincides with u on \(E_{1}\) . Then u is a Riesz endomorphism if and only if \(u_{1}\) is one;
\item
   Suppose that u is a Fredholm endomorphism of index 0. Let \(\widehat{\mathbf{E}}\) be the completion of E, and let \(\widehat{u} \in \mathcal{L}(\widehat{\mathbf{E}})\) be the continuous extension of u. Then u is a Riesz endomorphism if and only if \(\widehat{u}\) is one;
\item
 Endow the dual E' of E with the topology of bounded convergence, the topology of compact convergence, or the weak topology. If u is a Riesz endomorphism of E, then \(^{t}\) u is a Riesz endomorphism of E'.
\end{enumerate}
a) For \(u\) to be a Fredholm endomorphism of index zero, it is necessary and sufficient that \(u_1\) be one (III, p. 44, prop. 3). On the other hand, for every \(n \in \mathbf{N}\) , we have \(\mathrm{Ker}(u_1^n) = \mathrm{E}_1 \cap \mathrm{Ker}(u^n)\) , whence


\[
\dim \operatorname{Ker} (u _ {1} ^ {n}) \leqslant \dim \operatorname{Ker} (u ^ {n}) \leqslant \dim \operatorname{Ker} (u _ {1} ^ {n}) + \operatorname{codim} _ {\mathrm{E}} (\mathrm{E} _ {1}),
\]

therefore the sequence \((\dim\operatorname{Ker}(u^{n}))_{n\in\mathbf{N}}\) is bounded if and only if the sequence \((\dim\operatorname{Ker}(u_{1}^{n}))_{n\in\mathbf{N}}\) is. Consequently, u is a Riesz endomorphism of E if and only if \(u_{1}\) is a Riesz endomorphism of \(E_{1}\) (lemma 2).

\begin{enumerate}
\def\labelenumi{\alph{enumi})}
\setcounter{enumi}{1}
\item
  By prop. 4 of III, p. 44, for every \(n \in \mathbf{N}\) , the map \(\widehat{u}^n\) is a Fredholm endomorphism of \(\widehat{\mathrm{E}}\) such that \(\operatorname{Ker}(\widehat{u}^n) = \operatorname{Ker}(u^n)\) and \(\operatorname{ind}(\widehat{u}^n) = \operatorname{ind}(u^n) = n\operatorname{ind}(u)\) . Consequently, \(u\) is a Riesz endomorphism if and only if \(\widehat{u}\) is a Riesz endomorphism.
\item
  The transpose \({}^{t}u\) is a Fredholm endomorphism of index 0 of E' (III, p. 43, no. 3). Since the kernel of \((^{t}u)^{n}\) is the orthogonal complement of the image of \(u^{n}\) (EVT, IV, p. 27, prop. 2), the sequence \((\mathrm{Ker}({}^{t}u)^{n})_{n}\) is stationary. This proves that \({}^{t}u\) is a Riesz endomorphism of E'.
\end{enumerate}

PROPOSITION 8. --- Let E be a separated locally convex space and u an endomorphism of E. The following conditions are equivalent:

\begin{enumerate}
\def\labelenumi{(\roman{enumi})}
\item
  u is a Riesz endomorphism of E;
\item
  There exists a quasi-inverse v of u that commutes with u;
\item
  There exists a closed vector subspace \(E_{1}\) of E, of finite codimension, stable under u, such that u induces an automorphism of \(E_{1}\) .
\end{enumerate}

Implication (i) \(\Longrightarrow\) (ii) follows from prop. 6, d).

Suppose that \(u\) has a quasi-inverse \(v\) which commutes with \(u\) . The kernel \(\mathrm{E}_1\) of \(1_{\mathrm{E}} - u \circ v\) is a closed vector subspace of E, of finite codimension, stable under \(u\) and \(v\) . Since one has \(u(v(x)) = v(u(x)) = x\) for all \(x \in \mathrm{E}_1\) , the maps \(u\) and \(v\) induce automorphisms of \(\mathrm{E}_1\) inverses of each other. Hence (ii) implies (iii).

Finally, the implication (iii) \(\Longrightarrow\) (i) follows from the assertion \(a\) ) of Prop. 7 and from the fact that an automorphism is a Riesz endomorphism.

PROPOSITION 9. --- Let E be a separated locally convex space, and let u, v be permutable elements of \(\mathcal{L}(\mathrm{E})\) . The following conditions are equivalent:

\begin{enumerate}
\def\labelenumi{(\roman{enumi})}
\item
  The endomorphisms u and v are Riesz endomorphisms.
\item
  The endomorphism \(u \circ v\) is a Riesz endomorphism.
\end{enumerate}

Suppose that u and v are Riesz endomorphisms; denote by \(u'\) and \(v'\) their canonical quasi-inverses. The endomorphisms u, v, \(u'\) and \(v'\) of E commute (prop. 6, d)). The endomorphism \(v' \circ u'\) is a quasi-inverse of \(u \circ v\) which commutes with \(u \circ v\) , hence \(u \circ v\) is a Riesz endomorphism (prop. 8).

Suppose conversely that \(u \circ v\) is a Riesz endomorphism, and let w be its canonical quasi-inverse. Since u commutes with \(u \circ v\) , the endomorphisms u and w commute by prop. 6, d). Likewise v and w commute. Hence the endomorphisms u, v and w of E commute; it follows that the endomorphism \(v \circ w\) of E is a quasi-inverse of u and that \(w \circ u\) is a quasi-inverse of v. By prop. 8, u and v are Riesz endomorphisms.

\pandocbounded{\includegraphics[keepaspectratio]{images/part001_6cebf9a3.jpg}}

Let u be a Riesz endomorphism of E. If v is an automorphism of E, the endomorphism \(u \circ v\) of E is not necessarily a Riesz endomorphism of E, even if \(u \circ v - v \circ u\) has finite rank. If h is a finite-rank endomorphism of E, the endomorphism \(u + h\) is not necessarily a Riesz endomorphism (cf. III, p. 120, exercise 3).

PROPOSITION 10. --- Let E and F be locally convex spaces, \(p \in \mathcal{L}(\mathrm{E};\mathrm{F})\) and \(q \in \mathcal{L}(\mathrm{F};\mathrm{E})\) . Set \(u = 1_{\mathrm{E}} - q \circ p\) and \(v = 1_{\mathrm{F}} - p \circ q\) .

\begin{enumerate}
\def\labelenumi{\alph{enumi})}
\item
  Let \(n \in \mathbf{N}\) . The map \(p\) defines by restriction an isomorphism of topological vector spaces from \(\mathrm{Ker}(u^n)\) onto \(\mathrm{Ker}(v^n)\) and defines by passage to quotients an isomorphism of topological vector spaces from Coker( \(u^n\) ) onto Coker( \(v^n\) );
\item
  If the image of u is closed in E, that of v is closed in F;
\item
   If \(\operatorname{Ker}(u)\) has a topological supplement in E, then \(\operatorname{Ker}(v)\) has one in F. If \(\operatorname{Im}(u)\) has a topological supplement in E, then \(\operatorname{Im}(v)\) has one in F;
\item
  If u is a strict morphism, then v is a strict morphism.
\item
  If u is an automorphism of E, then v is an automorphism of F.
\item
  If \(u\) is a Fredholm endomorphism of E, then \(v\) is a Fredholm endomorphism of F, and we have \(\operatorname{ind}(u) = \operatorname{ind}(v)\) ;
\item
  Suppose E and F are separated. If u is a Riesz endomorphism of E, then v is a Riesz endomorphism of F.
\end{enumerate}

Let us first note the formulas

\[
q \circ v = u \circ q, \quad v \circ p = p \circ u.\tag{6}
\]

Let us then prove two lemmas.

Lemma 4. --- Let \(n \in \mathbf{N}\) . Set

\[
q _ {n} = \sum_ {i = 1} ^ {n} (- 1) ^ {i - 1} \binom {n} {i} q \circ (p \circ q) ^ {i - 1}.
\]

One has \(u^{n}=1_{E}-q_{n}\circ p\) and \(v^{n}=1_{F}-p\circ q_{n}\) .

One has \((q\circ p)^{i} = q\circ (p\circ q)^{i - 1}\circ p\) for every integer \(i\geqslant 1\) . We then compute

\[
u ^ {n} = (1 _ {\mathrm{E}} - q \circ p) ^ {n} = 1 _ {\mathrm{E}} - \sum_ {i = 1} ^ {n} (- 1) ^ {i - 1} {\binom {n} {i}} (q \circ p) ^ {i} = 1 _ {\mathrm{E}} - q _ {n} \circ p
\]

\[
v ^ {n} = (1 _ {\mathrm{F}} - p \circ q) ^ {n} = 1 _ {\mathrm{F}} - \sum_ {i = 1} ^ {n} (- 1) ^ {i - 1} \binom {n} {i} (p \circ q) ^ {i} = 1 _ {\mathrm{F}} - p \circ q _ {n},
\]

whence the result.

Lemma 5. --- Let \(u'\) be an element of \(\mathcal{L}(\mathrm{E})\) . Set \(v' = 1_{\mathrm{F}} + p \circ u' \circ q\) . If \(u'\) is an inverse of \(u\) , then \(v'\) is an inverse of \(v\) . If \(u'\) is a quasi-inverse of \(u\) , then \(v'\) is a quasi-inverse of \(v\) . If \(u'\) commutes with \(u\) , then \(v'\) commutes with \(v\) .

Using formulas (6), we compute

\[
v - v ^ {\prime} \circ v = - p \circ u ^ {\prime} \circ q \circ v = - p \circ u ^ {\prime} \circ u \circ q
\]

\[
v - v \circ v ^ {\prime} = - v \circ p \circ u ^ {\prime} \circ q = - p \circ u \circ u ^ {\prime} \circ q,
\]

whence

\[
1 _ {\mathrm{F}} - v ^ {\prime} \circ v = p \circ (1 _ {\mathrm{E}} - u ^ {\prime} \circ u) \circ q
\]

\[
1 _ {\mathrm{F}} - v \circ v ^ {\prime} = p \circ (1 _ {\mathrm{E}} - u \circ u ^ {\prime}) \circ q,
\]

which proves the lemma.

Given Lemma 4, it suffices to prove assertion \(a\) ) for \(n = 1\) . According to formulas (6), the maps \(p\) and \(q\) define, by passing to subspaces, continuous linear maps \(p'\) : \(\operatorname{Ker}(u) \to \operatorname{Ker}(v)\) and \(q'\) : \(\operatorname{Ker}(v) \to \operatorname{Ker}(u)\) , and by passing to quotients of continuous linear maps \(p''\) : \(\operatorname{Coker}(u) \to \operatorname{Coker}(v)\) and \(q''\) : \(\operatorname{Coker}(v) \to \operatorname{Coker}(u)\) .

Since \(1_{\mathrm{E}} - u = q \circ p\) and \(1_{\mathrm{F}} - v = p \circ q\) , the maps \(p'\) and \(q'\) on the one hand, \(p''\) and \(q''\) on the other hand, are mutually inverse isomorphisms. This proves \(a\) ).

For the image of a continuous linear map to be closed, it is necessary and sufficient that the cokernel of this map be a separated space. Thus b) follows from a).

Denote by \(i\colon \operatorname{Ker}(u)\to \mathrm{E}\) and \(j\colon \operatorname{Ker}(v)\to \mathrm{F}\) the canonical injections. Suppose that \(\operatorname{Ker}(u)\) has a topological supplement in E, and let \(r\colon \mathrm{E}\to \operatorname{Ker}(u)\) be a continuous linear retraction of \(i\) . Set \(r' = p'\circ r\circ q\) . We have \(r'\in \mathcal{L}(\mathrm{F};\operatorname {Ker}(v))\) and

\[
r ^ {\prime} \circ j = p ^ {\prime} \circ r \circ q \circ j = p ^ {\prime} \circ r \circ i \circ q ^ {\prime} = p ^ {\prime} \circ q ^ {\prime} = 1 _ {\operatorname{Ker} (v)},
\]

hence \(r'\) is a continuous linear retraction of j and \(\operatorname{Ker}(v)\) has a topological supplement in F. This proves the first assertion of c). The second is proved analogously, noting that if \(s: \operatorname{Coker}(u) \to \operatorname{E}\) is a continuous linear section of the canonical surjection from E onto \(\operatorname{Coker}(u)\) , then \(p \circ s \circ q''\) is a continuous linear section of the canonical surjection of F onto \(\operatorname{Coker}(v)\) .

Denote by \(\overline{u}\colon\mathrm{E}/\mathrm{Ker}(u)\to\mathrm{Im}(u)\) and \(\overline{v}\colon\mathrm{F}/\mathrm{Ker}(v)\to\mathrm{Im}(v)\) the continuous linear bijective mappings deduced from u and v. By formulas (6), the map p defines, by passing to quotients, a continuous linear map \(p_{1}\colon\mathrm{E}/\mathrm{Ker}(u)\to\mathrm{F}/\mathrm{Ker}(v)\) and p, q define, by passing to subspaces, continuous linear maps \(p_{0}\colon\mathrm{Im}(u)\to\mathrm{Im}(v)\) and \(q_{0}\colon\mathrm{Im}(v)\to\mathrm{Im}(u)\) .

Denote by \(t \colon \operatorname{Im}(v) \to \operatorname{F} / \operatorname{Ker}(v)\) the composite map of the canonical injection \(\operatorname{Im}(v) \to \operatorname{F}\) and of the canonical surjection \(\operatorname{F} \to \operatorname{F} / \operatorname{Ker}(v)\) ; it is continuous. We have \(\overline{v} \circ t = 1_{\operatorname{Im}(v)} - p_0 \circ q_0\) and \(\overline{v} \circ p_1 = p_0 \circ \overline{u}\) , whence

\[
\overline {{v}} \circ (t + p _ {1} \circ \overline {{u}} ^ {- 1} \circ q _ {0}) = \overline {{v}} \circ t + p _ {0} \circ q _ {0} = 1 _ {\mathrm{Im} (v)}.
\]

This proves that \(t + p_{1} \circ \overline{u}^{-1} \circ q_{0}\) is the inverse bijection of \(\overline{v}\) . If u is strict, the map \(\overline{u}^{-1}\) is continuous, and so is \(\overline{v}^{-1}\) which proves that v is strict. This proves d).

By Lemma 5, if \(u\) is an automorphism of E, then \(v\) is an automorphism of F. Similarly, if \(u\) is a Fredholm endomorphism of E, then \(v\) is a Fredholm endomorphism of F, and one then has \(\operatorname{ind}(u) = \operatorname{ind}(v)\) according to \(a\) ). This proves \(e\) ) and \(f\) ). Finally, if E and F are separated and \(u\) is a Riesz endomorphism of E, then \(v\) is a Riesz endomorphism of F by Lemma 5 and condition (ii) of Prop. 8; this is assertion \(g\) ).

\subsection{Fredholm mappings and Riesz mappings between Fréchet spaces}

PROPOSITION 11. --- Let E and F be Fréchet spaces. For a continuous linear map u from E into F to be a Fredholm map, it is necessary and sufficient that its kernel and cokernel be of finite dimension.

This follows from the characterization of Fredholm maps given by condition (ii) of Prop. 2 in III, p. 42, and from the following lemma:

Lemma 6. --- Let E and F be Fréchet spaces and \(u \in \mathcal{L}(\mathrm{E};\mathrm{F})\) . We assume that the image of \(u\) has finite codimension in F. Then the image of \(u\) is closed in F and \(u\) is a strict morphism.

Let G be a vector subspace of F supplementary to \(\operatorname{Im}(u)\) ; the subspace G is closed (EVT, I, p. 14, cor. 1), hence the quotient space F/G is a Fréchet space (TG, IX, p. 25, prop. 4). Let \(\pi: F \to F/G\) be the canonical surjection. The map \(\pi \circ u\) is surjective, hence strict (EVT, I, p. 19, cor. 3). The assertion then follows from cor. 2 of III, p. 40.

PROPOSITION 12. --- Let E and F be Fréchet spaces. Endow their duals E' and F' with the weak topology, the topology of compact convergence, or the topology of bounded convergence. Let \(u \in \mathcal{L}(\mathrm{E};\mathrm{F})\) . For u to be a Fredholm map from E to F, it is necessary and sufficient that \({}^{t}u\) be a Fredholm map from F' to E'.

The condition is necessary (III, p. 43, no. 3). Let us prove that it is sufficient. If \(^{t}u\) is a Fredholm map, it is a strict morphism (III, p. 42, prop. 2), hence u is a strict morphism with closed image (EVT, IV, p. 28, th. 1 and p. 29, cor. 3). Taking into account the canonical isomorphisms Coker( \(^{t}u\) ) → Ker(u)' and Ker( \(^{t}u\) ) → Coker(u)' (EVT, IV, p. 27, prop. 2), the vector spaces Ker(u) and Coker(u) are finite-dimensional and u is a Fredholm operator (prop. 2 of III, p. 42).

PROPOSITION 13. --- Let E be a Fréchet space and \(u \in \mathcal{L}(\mathrm{E})\) .

\begin{enumerate}
\def\labelenumi{\alph{enumi})}
\item
  The endomorphism u is a Riesz endomorphism if and only if its nilespace N is finite-dimensional and its conilespace I is of finite codimension;
\item
 Endow E' with the topology of bounded convergence, the topology of compact convergence, or the weak topology. If \(^{t}\) u is a Riesz endomorphism of E', then u is a Riesz endomorphism of E.
\end{enumerate}
a) Suppose that the nilespace N is finite-dimensional and that the conilespace I is of finite codimension. Since \(\operatorname{Ker}(u) \subset \operatorname{N}\) and \(\operatorname{I} \subset \operatorname{Im}(u)\) , the kernel of u is finite-dimensional and the image of u is of finite codimension. By the definition and proposition 11, the map u is a Riesz endomorphism. The converse follows from the definition.

b) By proposition 12, the hypothesis implies that \(u\) is a Fredholm endomorphism of index \(\text{ind}(^t u) = -\text{ind}(u) = 0\) ( \(n^o 3\) , formula (4)). Let \(n \in \mathbf{N}\) . Since the image of \(u^n\) is closed and has as its orthogonal the kernel of \(^t u^n\) (EVT, IV, p. 27, prop. 2), the sequence \((\text{Im}(u^n))\) is stationary, and \(u\) is therefore a Riesz endomorphism.


\subsection{Spectral characterization of Riesz endomorphisms}

Let E be a complex Banach space and u an endomorphism of E. Recall that we denote by \(\mathrm{Sp}(u)\) the spectrum of u relative to the unital Banach algebra \(\mathcal{L}(\mathrm{E})\) (cf. § 7 of I, p. 127).

Suppose that 0 is an isolated point of \(\mathrm{Sp}(u)\) ; recall that we then denote by \(e_0(u)\) the spectral projector associated with \(u\) and with the open and closed subset \(\{0\}\) of the spectrum of \(u\) (cf. n°3 of I, p. 131). We have \(e_0(u) = f(u)\) for every germ of holomorphic function \(f\) in the neighborhood of \(\mathrm{Sp}(u)\) which equals 1 in a neighborhood of 0 and vanishes in a neighborhood of \(\mathrm{Sp}(u) \setminus \{0\}\) .

By passing to subspaces, the endomorphism \(u\) induces a quasi-nilpotent endomorphism of the image of \(e_0(u)\) , whose spectrum is reduced to \(\{0\}\) , and an automorphism of the kernel of \(e_0(u)\) , whose spectrum is \(\mathrm{Sp}(u) \setminus \{0\}\) (loc. cit.).

PROPOSITION 14. --- Let E be a complex Banach space. For an element u of \(\mathcal{L}(\mathrm{E})\) to be a Riesz endomorphism of E, it is necessary and sufficient that one of the following two mutually exclusive conditions be satisfied:

\begin{enumerate}
\def\labelenumi{(\roman{enumi})}
\item
  The endomorphism u of E is an automorphism of E;
\item
  The point 0 is an isolated point of \(\operatorname{Sp}(u)\) and the projector \(e_0(u)\) has finite rank.
\end{enumerate}

When condition (ii) is satisfied, the image of \(e_{0}(u)\) is the nilespace of u and the kernel of \(e_{0}(u)\) is the conilespace of u.

Every automorphism of E is a Riesz endomorphism of E (III, p. 47, remark 1). If u satisfies condition (ii), the kernel F of \(e_{0}(u)\) is a closed vector subspace of finite codimension of E, stable under u. Let \(u_{F}\) be the endomorphism of F deduced from u. Its spectrum is contained in \(\mathbf{C}\setminus\{0\}\) (n°3 of I, p. 131), hence \(u_{F}\) is an automorphism of F. It follows that u is a Riesz endomorphism of E (III, p. 48, prop. 8).

Conversely, let \(u\) be a Riesz endomorphism of E. Denote by N its nilespace and by I its conilespace. By prop. 6 of III, p. 47, the space E is the topological direct sum of N and I, and the map \(u\) defines by passage to subspaces a nilpotent endomorphism \(u_{\mathrm{N}}\) of N and an automorphism \(u_{\mathrm{I}}\) of I. In particular, we have \(\mathrm{Sp}(u_{\mathrm{N}}) \subset \{0\}\) and \(0 \not\in \mathrm{Sp}(u_{\mathrm{I}})\) . If N is zero, then I = E and \(u\) is an automorphism of E. Otherwise, 0 is an isolated point of \(\mathrm{Sp}(u)\) and \(e_0(u)\) is the projector with image N and kernel I (prop. 2 of I, p. 129); in particular, it is of finite rank.

Remark. --- Let E be a real Banach space and let u be an endomorphism of E. Suppose that 0 is an isolated point of the complex spectrum of u, that is, of the spectrum of the endomorphism \(1 \otimes u\) of the complete normable complexified space \(\mathrm{E}_{(\mathbf{C})}\) of E (I, p. 85, n° 13). The subset \(\{0\}\) of the complex spectrum of \(u\) is open and closed, and invariant under conjugation; denote by \(e_0(u) \in \mathcal{L}(\mathrm{E})\) the spectral projector associated with it (n° 13 of I, p. 85). Proposition 14 holds mutatis mutandis in this setting.

\section{PERTURBATIONS IN BANACH SPACES}

\subsection{Direct morphisms}

DEFINITION 1. --- Let E be a topological vector space. A vector subspace of E is said to be direct if it admits a topological complement.

For a vector subspace F of E to be direct, it is necessary and sufficient that there exist a projector of image F in \(\mathcal{L}(\mathrm{E})\) . If \(p\) is such a projector, E is the topological direct sum of F and the kernel of \(p\) . The kernel of \(1_{\mathrm{E}} - p\) is F; if E is separated, then F is therefore closed.

Let F be a direct vector subspace of E and \(E_{1}\) a vector subspace of E containing F; then F is a direct vector subspace of \(E_{1}\) .

Every closed vector subspace of finite codimension in E is a direct summand (EVT, I, p. 15, prop. 3), and every vector subspace F of E contained in the closure of \{0\} is direct (indeed, there exists a projector with image F, and such a projector is necessarily continuous).

PROPOSITION 1. --- Let E be a locally convex space. Every finite-dimensional vector subspace of E is complemented.

Let F be a finite-dimensional vector subspace of E. Let S denote a topological complement of the closure of \(\{0\}\) in F, which is a direct subspace of E. The space \(F_{2} = F \cap S\) is separated and of finite dimension; there therefore exists an integer \(n \geqslant 0\) such that \(F_{2}\) is isomorphic to \(K^{n}\) (EVT, I, p. 14, th. 2). There then exists a continuous linear map \(p: S \to F_{2}\) extending the identity map of \(F_{2}\) . The kernel of p is a topological supplement of \(F_{2}\) in S, and also of F in E.

DEFINITION 2. --- Let E and F be locally convex spaces, and let u be a continuous linear map from E into F. One says that u is a direct morphism if u is a strict morphism whose kernel is a direct vector subspace of E and whose image is a direct vector subspace of F.

Let \(u \in \mathcal{L}(\mathrm{E};\mathrm{F})\) . For the kernel and image of \(u\) to be direct vector subspaces of E and F, respectively, it is necessary and sufficient that there exist decompositions \(\mathrm{E} = \mathrm{E}_1 \oplus \mathrm{E}_2\) and \(\mathrm{F} = \mathrm{F}_1 \oplus \mathrm{F}_2\) into topological direct sums such that \(u\) be represented by a matrix \(\begin{pmatrix} u_1 & 0 \\ 0 & 0 \end{pmatrix}\) where \(u_1 \in \mathcal{L}(\mathrm{E}_1; \mathrm{F}_1)\) is bijective. Since the kernel of \(u\) is then \(\mathrm{E}_2\), the canonical linear map from \(\mathrm{E}_1\) onto \(\mathrm{E}/\mathrm{E}_2\) is an isomorphism, and the image of \(u\) is \(\mathrm{F}_1\), one sees that \(u\) is a strict morphism if and only if \(u_1\) is an isomorphism of \(\mathrm{E}_1\) onto \(\mathrm{F}_1\). Denote by \(v\) the element of \(\mathcal{L}(\mathrm{F};\mathrm{E})\) represented by the matrix \(\begin{pmatrix} u_1^{-1} & 0 \\ 0 & 0 \end{pmatrix}\). Then the linear map \(u \circ v\) is the projector in F with kernel \(\mathrm{F}_2\) and image \(\mathrm{F}_1\), the linear map \(v \circ u\) is the projector in E with kernel \(\mathrm{E}_2\) and image \(\mathrm{E}_1\), and one has \(u \circ v \circ u = u\).

Conversely, one has the following result:

PROPOSITION 2. --- Let E and F be locally convex spaces, let \(u \in \mathcal{L}(\mathrm{E};\mathrm{F})\) and \(v \in \mathcal{L}(\mathrm{F};\mathrm{E})\) . Suppose that \(u = u \circ v \circ u\) . Then \(u\) is a direct morphism. Moreover, \(v \circ u\) is a continuous projector in E with kernel \(\operatorname{Ker}(u)\) and \(u \circ v\) is a continuous projector in F with image \(\operatorname{Im}(u)\) .

Set \(p = v \circ u\) and \(q = u \circ v\) . We have

\[
p ^ {2} = v \circ (u \circ v \circ u) = v \circ u = p, q ^ {2} = (u \circ v \circ u) \circ v = u \circ v = q,
\]

hence p and q are projectors. They are continuous. Let \(E_{1}\) and \(E_{2}\) be the image and the kernel of p, and let \(F_{1}\) and \(F_{2}\) be the image and the kernel of q. Since \(u \circ v \circ u = u\) , we have

\[
\operatorname{Ker} (u) \subset \operatorname{Ker} (v \circ u) \subset \operatorname{Ker} (u \circ v \circ u) = \operatorname{Ker} (u),
\]

hence \(\operatorname{Ker}(u) = \operatorname{Ker}(v \circ u) = \operatorname{Ker}(p) = \operatorname{E}_{2}\) . Similarly, one deduces from the inclusions

\[
\operatorname{Im} (u) \supset \operatorname{Im} (u \circ v) \supset \operatorname{Im} (u \circ v \circ u) = \operatorname{Im} (u)
\]

that the image of \(u\) is \(\mathrm{F}_{1}\) . The spaces \(\mathrm{E}_{2}\) and \(\mathrm{F}_{1}\) are direct vector subspaces of E and F respectively.

The vector space E is the topological direct sum of \(E_{1}\) and \(E_{2}\) , and u defines by restriction a bijective linear map \(u_{1}\) of \(E_{1}\) onto \(F_{1}\) . For every \(x \in E_{1}\) , we have \(v(u(x)) = p(x) = x\) , and the map \(u_{1}^{-1}\) of \(F_{1}\) onto \(E_{1}\) coincides with v on \(F_{1}\) . Consequently, \(u_{1}\) is an isomorphism of \(E_{1}\) onto \(F_{1}\) and u is a strict morphism from E into F.

Let E and F be locally convex spaces and \(u \in \mathcal{L}(\mathrm{E};\mathrm{F})\) . For u to be an injective (resp. surjective) direct morphism, it is necessary and sufficient that there exist a \(v\) in \(\mathcal{L}(\mathrm{F};\mathrm{E})\) such that \(v \circ u = 1_{E}\) (resp. \(u \circ v = 1_{F}\) ) (cf. TG, III, p. 47 and 48).

PROPOSITION 3. --- Let E and F be Banach spaces. The following sets are open subsets of the Banach space \(\mathcal{L}(\mathrm{E};\mathrm{F})\) :

\begin{enumerate}
\def\labelenumi{\alph{enumi})}
\item
  The set \(\mathcal{I}(\mathrm{E};\mathrm{F})\) of isomorphisms from E onto F;
\item
   The set \(\mathcal{M}\mathcal{D}(\mathrm{E};\mathrm{F})\) of injective direct morphisms from E into F, and more precisely, for every closed vector subspace \(F_{1}\) of F, the set \(\mathcal{M}_{\mathrm{F}_{1}}(\mathrm{E};\mathrm{F})\) of elements of \(\mathcal{M}\mathcal{D}(\mathrm{E};\mathrm{F})\) whose image is a topological supplement of \(F_{1}\) ;
\item
   The set \(\mathcal{E}\mathcal{D}(\mathrm{E};\mathrm{F})\) of surjective direct morphisms from E into F, and more precisely, for every closed vector subspace \(\mathrm{E}_1\) of E, the set \(\mathcal{E}_{\mathrm{E}_1}(\mathrm{E};\mathrm{F})\) of elements of \(\mathcal{E}\mathcal{D}(\mathrm{E};\mathrm{F})\) whose kernel is a topological supplement of \(\mathrm{E}_1\) .
\end{enumerate}

Moreover, the map \(u \mapsto u^{-1}\) of \(\mathcal{I}(\mathrm{E};\mathrm{F})\) onto \(\mathcal{I}(\mathrm{F};\mathrm{E})\) is analytic.

By definition, \(\mathcal{I}(\mathrm{E};\mathrm{E})\) is the set of invertible elements of the complete normed algebra \(\mathcal{L}(\mathrm{E})\) . By TG, IX, p. 40, prop. 14, it is an open subset of \(\mathcal{L}(\mathrm{E})\) . The map \(u\mapsto u^{-1}\) from \(\mathcal{I}(\mathrm{E};\mathrm{E})\) into \(\mathcal{I}(\mathrm{E};\mathrm{E})\) is analytic (LIE, III, § 1, n° 1, example 2).

If the set \(\mathcal{I}(\mathrm{E};\mathrm{F})\) is empty, it is open. Otherwise, let \(u_0\) be an isomorphism from E onto F. The map \(v\mapsto u_0\circ v\) is then an isomorphism of \(\mathcal{L}(\mathrm{E})\) onto \(\mathcal{L}(\mathrm{E};\mathrm{F})\) which transforms \(\mathcal{I}(\mathrm{E};\mathrm{E})\) into \(\mathcal{I}(\mathrm{E};\mathrm{F})\) . The set \(\mathcal{I}(\mathrm{E};\mathrm{F})\) is therefore open in \(\mathcal{L}(\mathrm{E};\mathrm{F})\) . Moreover, if \(u = u_0\circ v\) is an element of \(\mathcal{I}(\mathrm{E};\mathrm{F})\) , we have \(u^{-1} = v^{-1}\circ u_0^{-1}\) , and the map \(u\mapsto u^{-1}\) of \(\mathcal{I}(\mathrm{E};\mathrm{F})\) onto \(\mathcal{I}(\mathrm{F};\mathrm{E})\) is therefore analytic.

Let \(F_{1}\) be a closed vector subspace of F. For every \(u \in \mathcal{L}(E;F)\) , let \(\overline{u}\) be the element of \(\mathcal{L}(E \times F_{1};F)\) defined by \(\overline{u}(x,y) = u(x) + y\) . The map \(u \mapsto \overline{u}\) from \(\mathcal{L}(E;F)\) into \(\mathcal{L}(E \times F_{1};F)\) is continuous, and \(\mathcal{M}_{\mathrm{F}_{1}}(E;F)\) is the set of elements u of \(\mathcal{L}(E;F)\) such that \(\overline{u}\) belongs to \(\mathcal{I}(E \times F_{1};F)\) . Since \(\mathcal{I}(E \times F_{1};F)\) is open in \(\mathcal{L}(E \times F_{1};F)\)

by what precedes, the set \(\mathcal{M}_{\mathrm{F}_1}(\mathrm{E};\mathrm{F})\) is open in \(\mathcal{L}(\mathrm{E};\mathrm{F})\) . The same is true of \(\mathcal{MD}(\mathrm{E};\mathrm{F})\) , which is the union of the \(\mathcal{M}_{\mathrm{F}_1}(\mathrm{E};\mathrm{F})\) when \(\mathrm{F}_1\) ranges over the set of closed vector subspaces of F.

Let \(E_{1}\) be a closed vector subspace of E, and let p be the canonical map of E onto the quotient Banach space \(E/E_{1}\) . For an element \(u\) of \(\mathcal{L}(E;F)\) to belong to \(\mathcal{E}_{\mathrm{E}_{1}}(\mathrm{E};\mathrm{F})\), it is necessary and sufficient that the map \((u,p)\) from E into \(F\times E/E_{1}\) belong to \(\mathcal{I}(E;F\times E/E_{1})\). As before, we deduce that \(\mathcal{E}_{\mathrm{E}_{1}}(\mathrm{E};\mathrm{F})\) is open in \(\mathcal{L}(\mathrm{E};\mathrm{F})\); the same is true of \(\mathcal{E}\mathcal{D}(\mathrm{E};\mathrm{F})\), which is the union of the \(\mathcal{E}_{\mathrm{E}_{1}}(\mathrm{E};\mathrm{F})\).

\subsection{Perturbation of Fredholm maps}

THEOREM 1. --- Let E and F be Banach spaces. The set \(\mathcal{F}(\mathrm{E};\mathrm{F})\) of Fredholm operators from E to F is open in the Banach space \(\mathcal{L}(\mathrm{E};\mathrm{F})\) , and the map \(u\mapsto\mathrm{ind}(u)\) of \(\mathcal{F}(\mathrm{E};\mathrm{F})\) into \(\mathbf{Z}\) is locally constant.

The theorem follows from the following more precise statement:

PROPOSITION 4. --- Let E and F be Banach spaces, \(u_{0}\colon\mathrm{E}\to\mathrm{F}\) a Fredholm operator and \(v_{0}\) a quasi-inverse of \(u_{0}\) . There exists an open neighborhood U of \(u_{0}\) in \(\mathcal{L}(\mathrm{E};\mathrm{F})\) and an analytic mapping \(\varphi\colon\mathrm{U}\to\mathcal{L}(\mathrm{F};\mathrm{E})\) such that

\begin{enumerate}
\def\labelenumi{(\roman{enumi})}
\item
  One has \(\varphi(u_{0}) = v_{0}\) ;
\item
  For every \(u\) in \(U\) , the map \(\varphi(u)\) is a quasi-inverse of \(u\) , and in particular \(u\) is a Fredholm map;
\item
  For every \(u\) in U, we have \(\operatorname{ind}(u) = \operatorname{ind}(u_0)\).
\end{enumerate}

By prop. 2 of III, p. 42, (ii), there exist topological direct sum decompositions \(\mathrm{E} = \mathrm{E}_1 \oplus \mathrm{E}_2\) and \(\mathrm{F} = \mathrm{F}_1 \oplus \mathrm{F}_2\) , and there exists \(\alpha_0 \in \mathcal{I}(\mathrm{E}_1; \mathrm{F}_1)\) such that \(\mathrm{E}_2\) and \(\mathrm{F}_2\) are finite-dimensional and \(u_0\) is represented by the matrix \(\left( \begin{array}{cc}\alpha_0 & 0\\ 0 & 0 \end{array} \right)\) relative to these decompositions.

Let U be the set of elements \(u\) of \(\mathcal{L}(\mathrm{E};\mathrm{F})\) such that, in the matrix representation \(\left( \begin{array}{cc}\alpha & \beta \\ \gamma & \delta \end{array} \right)\) of \(u\) with respect to these decompositions, we have \(\alpha \in \mathcal{I}(\mathrm{E}_1;\mathrm{F}_1)\) . Since \(\mathcal{I}(\mathrm{E}_1;\mathrm{F}_1)\) is open in \(\mathcal{L}(\mathrm{E}_1;\mathrm{F}_1)\)

(prop. 3 of III, p. 57), U is an open neighborhood of \(u_0\) in \(\mathcal{L}(\mathrm{E};\mathrm{F})\) . For \(u = \begin{pmatrix} \alpha & \beta \\ \gamma & \delta \end{pmatrix}\) in U, set

\[
\varphi (u) = v _ {0} + \left( \begin{array}{c c} \alpha^ {- 1} - \alpha_ {0} ^ {- 1} & 0 \\ 0 & 0 \end{array} \right).\tag{1}
\]

One has \(\varphi(u_{0}) = v_{0}\) and the map \(\varphi\) is analytic (loc. cit.). Modulo continuous linear maps of finite rank, we have the congruences

\[
u \equiv \left( \begin{array}{c c} \alpha & 0 \\ 0 & 0 \end{array} \right) , v _ {0} \equiv \left( \begin{array}{c c} \alpha_ {0} ^ {- 1} & 0 \\ 0 & 0 \end{array} \right) , \varphi (u) \equiv \left( \begin{array}{c c} \alpha^ {- 1} & 0 \\ 0 & 0 \end{array} \right).
\]

Therefore, \(\varphi(u)\) is a quasi-inverse of \(u\) . Every element \(u\) from U defines by restriction an isomorphism from \(E _1\) onto a topological complement of \(F _2\) in F. By prop. 3 of III, p. 44, we therefore have

\[
\begin{array}{c} \operatorname{ind} (u) = \operatorname{codim} _ {\mathrm{F}} (u (\mathrm{E} _ {1})) - \operatorname{codim} _ {\mathrm{E}} (\mathrm{E} _ {1}) \\ = \dim (\mathrm{F} _ {2}) - \dim (\mathrm{E} _ {2}) = \operatorname{ind} (u _ {0}), \end{array}
\]

which concludes the proof.